% \iffalse meta-comment
%
% Copyright (C) \the\year by Liu Benyuan <liubenyuan@gmail.com>
% This file may be distributed and/or modified under the
% conditions of the LaTeX Project Public License, either
% version 1.2 of this license or (at your option) any later
% version. The latest version of this license is in:
%
% http://www.latex-project.org/lppl.txt
%
% and version 1.2 or later is part of all distributions of
% LaTeX version 1999/12/01 or later.
%
% \fi
% 
% \iffalse
% <package>\NeedsTeXFormat{LaTeX2e}[1999/12/01]
% <package>\ProvidesPackage{nudtpaper}
% <package>		[2009/08/12 v1.0 By Liu Benyuan <liubenyuan@gmail.com>]
%<*driver>
\ProvidesFile{nudtpaper.dtx}[2009/08/12 v1.0 NUDT]
\documentclass[11pt]{ltxdoc}
\usepackage{nudtx}
\EnableCrossrefs
\CodelineIndex
\RecordChanges
\begin{document}
  \DocInput{\jobname.dtx}
\end{document}
%</driver>
% \fi
% 
% \def\thuthesis{\textsc{Thu}\-\textsc{Thesis}}
% \def\nudtpaper{\textsc{Nudt}\-\textsc{Paper}}
% 
% \CheckSum{1224}
% \CharacterTable
%  {Upper-case    \A\B\C\D\E\F\G\H\I\J\K\L\M\N\O\P\Q\R\S\T\U\V\W\X\Y\Z
%   Lower-case    \a\b\c\d\e\f\g\h\i\j\k\l\m\n\o\p\q\r\s\t\u\v\w\x\y\z
%   Digits        \0\1\2\3\4\5\6\7\8\9
%   Exclamation   \!     Double quote  \"     Hash (number) \#
%   Dollar        \$     Percent       \%     Ampersand     \&
%   Acute accent  \'     Left paren    \(     Right paren   \)
%   Asterisk      \*     Plus          \+     Comma         \,
%   Minus         \-     Point         \.     Solidus       \/
%   Colon         \:     Semicolon     \;     Less than     \<
%   Equals        \=     Greater than  \>     Question mark \?
%   Commercial at \@     Left bracket  \[     Backslash     \\
%   Right bracket \]     Circumflex    \^     Underscore    \_
%   Grave accent  \`     Left brace    \{     Vertical bar  \|
%   Right brace   \}     Tilde         \~}
%
% \changes{v0.99}{2009/08/12}{Initial Release}
%
% \GetFileInfo{\jobname.dtx}
% 
% \DoNotIndex{\begin,\end,\begingroup,\endgroup}
% \DoNotIndex{\ifx,\ifdim,\ifnum,\ifcase,\else,\or,\fi}
% \DoNotIndex{\let,\def,\xdef,\newcommand,\renewcommand}
% \DoNotIndex{\expandafter,\csname,\endcsname,\relax,\protect}
% \DoNotIndex{\Huge,\huge,\LARGE,\Large,\large,\normalsize}
% \DoNotIndex{\small,\footnotesize,\scriptsize,\tiny}
% \DoNotIndex{\normalfont,\bfseries,\slshape,\interlinepenalty}
% \DoNotIndex{\hfil,\par,\hskip,\vskip,\vspace,\quad}
% \DoNotIndex{\centering,\raggedright}
% \DoNotIndex{\c@secnumdepth,\@startsection,\@setfontsize}
% \DoNotIndex{\ ,\@plus,\@minus,\p@,\z@,\@m,\@M,\@ne,\m@ne}
% \DoNotIndex{\@@par,\DeclareOperation,\RequirePackage,\LoadClass}
% \DoNotIndex{\AtBeginDocument,\AtEndDocument}
%
% \IndexPrologue{\section*{索引}%
%    \addcontentsline{toc}{section}{索~~~~引}}
% \GlossaryPrologue{\section*{修改记录}%
%    \addcontentsline{toc}{section}{修改记录}}
%
% \renewcommand{\abstractname}{摘~~要}
% \renewcommand{\contentsname}{目~~录}
%
% \title{\textsc{NUDTpaper:}\,\,NUDT研究生学位论文\LaTeX{}模板使用手册\thanks{NUDT \LaTeX{} Thesis Template}}
% \author{刘本源 \\ \texttt{Liubenyuan@gmail.com}}
% \date{\fileversion\ (\filedate)}
%
% \maketitle
% \thispagestyle{empty}
%
% \begin{abstract}
% 本模板旨在提供规范的国防科学技术大学\LaTeX{}写作模板环境，
% 现支持硕士/博士学位论文格式，可以自动生成盲评、制作A3封面。
% \end{abstract}
%
% \vspace{2cm}
% \def\abstractname{免责声明}
% \begin{abstract}\noindent
% \begin{enumerate}
% \item 本模板的发布遵守 \LaTeX{} Project Public License，使用前请认真阅读协议内容
% \item 本模板创立参照官方严格的论文写作手册，并同时参照硕士/博士学位论文\textbf{doc}文档对比修改
% \item 国防科学技术大学对论文写作提供写作指南与官方\textbf{doc}模板，
% 同时提供官方的\LaTeX{}模板，本模板的出发点是方便大家使用专业的高效的论文书写工具，
% 其有点在于注重排版质量、命令规范、使用方便、更新及时，符合论文撰写说明。
% 但任何由于使用本模板而引起的论文格式审查问题均与本模板作者无关。
% \item 任何个人或组织均可以本模板为基础进行修改、扩展，生成新的专用模板，但请严格遵
% 守\LaTeX{} Project Public License 协议
% \item 欢迎提出修改意见
% \end{enumerate}
% \end{abstract}
%
% \clearpage
% \tableofcontents
%
% \clearpage
% \pagenumbering{arabic}
% \pagestyle{mainpage}
%
% \section{快速上手}
% 这部分是专门为那些想快速开始写论文的人准备的。
% \begin{description}
% \item[安装\TeX] 下载最新的\TeX{}live或者C\TeX{}并安装
% \item[字体] 用户需要具备\verb|simsun.ttf|, \verb|simhei.ttf|, \verb|simkai.ttf|,
% \verb|STZHONGS.TTF|, 上述字体都是windows自带的; 除此之外，在网上搜索(或者C\TeX{}
% 论坛)``Adobe Opentype 中文字体''，一搜一大把，确保下载下来Adobe的四款OTF字体:
% 宋，黑，仿宋，楷体。Linux用户可将上述字体复制到\verb|/usr/share/fonts/TTF|下。
% \item[试一试] 解压缩下载的模板，双击makepdf.bat(祈祷一下)，如果生成了
% \verb|thesis.pdf|$\rightarrow$
% \item[那么我的那些常用的包都在么？] 你会想我的{\bf Trans}论文可以无缝
% 的复制过来么？ 对于这一点，你可以修改\verb|mynudt.sty|来实现。但是{\hei 注意}，大部分包
% 都在模板中了，而且{\hei 切记切记}，不要擅自改动字体等版面设计，我们继续看$\rightarrow$
% \item[咦，数学公式不是很美观呀] 笔者{\hei 强烈}建议用户使用{\bf mtpro2}宏包的，怎么使用，
% 又有哪些好处，参见bookzh.sty吧！不会错的。好了，我们专注于内容本身吧$\rightarrow$
% \item[开始写了] 所有文件均采用UTF8编码，因此要保证你的\TeX{}编辑器
% (winedt, texworks, texmaker, vim, 记事本($\cdots{}$)等)支持这种编码，
% (经过一番搜索设置后)打开\verb|thesis.tex|，如果看到的是中文$\rightarrow$
% \item[漫长的写作] 手边准备着\LaTeX{}的常用帮助文档(数学，图表，引用等),
% 结合你喜欢的文献管理软件(JabRef等), 漫长的\texttt{编辑，编译，修改，编辑，
% 编译$\cdots$}过程之后，突然有一天发现你写完了$\rightarrow$
% \item[校订] 经过老师师兄师弟师妹齐心协力校正之后，你所做的只是:
% \texttt{生成明评论文，制作明评封面，生成盲评论文，制作盲评封面}，
% 装订，上交$\rightarrow$
% \end{description}
% {\color{magenta} Done!}
%
% \section{模板介绍}
%
% \textsc{NUDTpaper} 旨在帮助并且推广\LaTeX{}在国防科技大学论文中的应用，
% 本文将尽可能帮助用户掌握\textsc{NUDTpaper}的安装方法，
% 如果仍旧有不清晰的地方可以参考样例文件或者
% 给作者邮件\footnote{liubenyuan@gmail.com}，感兴趣的同学可以帮忙维护模板，
% 这个模板首先符合官方的设计要求，希望同学们在使用后能够提出你们的修改意见。
% 该模板很大程度上参考了6院黄老师sofoot的国防科大博士论文模板，
% 哈工大的\LaTeX{}模板以及清华的Thuthesis
% \footnote{主页：\url{http://thuthesis.sourceforge.net}}，
% 有很多使用的帮助、\verb|.cls|中的命令以及版面设置均来自Thuthesis和sofoot的模板，
% 对此的引用表示感谢。
%
% {\color{blue}\fs 模板的作用在于减轻论文写作过程中格式调整的时间，
% 其前提就是遵守模板的用法，不提倡手动更改格式，不建议正文中使用
% 手动调节版面的命令，尤其禁止修改行距和使用\verb|\normalsize|，
% 否则即使使用了\textsc{NUDTpaper}也难以保证输出的论文符合学校规范。}
%
% \section{安装}
% \label{sec:install}
%
% \subsection{下载}
% \textsc{NUDTpaper} 主页：\url{http://nudtpaper.googlecode.com}。
% 模板的更新信息发布在\href{http://bbs.ctex.org}{Ctex论坛}。
% \nudtpaper{}的开发版本同样可以在\textsc{gitorious}上获得。
%
% \subsection{模板的组成部分}
% 下表列出了 \nudtpaper{} 的主要文件及其功能介绍，学习模板的最好办法
% 就是参考thesis.pdf！
% \begin{center}
% \begin{longtable}{l|p{8cm}}
% \toprule
% {\hei 文件（夹）} & {\hei 功能描述}\\\midrule
% \endfirsthead
% \toprule
% {\hei 文件（夹）} & {\hei 功能描述}\\\midrule
% \endhead
% \endfoot
% \endlastfoot
% nudtpaper.ins & 模板驱动文件 \\
% nudtpaper.dtx & 模板文档代码的混合文件\\
% nudtpaper.cls & 模板类文件\\
% nudtpaper.cfg & 模板配置文件\\
% thesis.bib & 参考文献样式文件\\
% \hline
% mynudt.sty & 在这里添加你自己的宏包 \\
% thesis.tex & 示例文档主文件\\
% ref/ & 示例文档参考文献目录\\
% data/ & 示例文档章节具体内容\\
% figures/ & 示例文档图片路径\\
% \textbf{nudtpaper.pdf} & 用户手册（本文档）\\
% \textbf{thesis.pdf} & 示例文档 \\
% \bottomrule
% \end{longtable}
% \end{center}
%
% \subsection{\TeX{}系统的选择}
% 有网络环境的用户推荐安装\href{http://www.tug.org/texlive}{\TeX{}live}，
% \href{http://miktex.org}{MiKTeX}或者\href{http://www.ctex.org}{C\TeX}，
% 对于无网络环境的，主要是针对教研室用户，推荐{\TeX{}live}或者C\TeX{}完整版，安装
% 过程很简单，一路下一步即可，但是需要\textbf{注意：}
%
% \begin{description}
% \item[字体] TTF选项默认调用Windows系统字体，其中楷体、仿宋需要安装Office；OTF选项需要
% Adobe的商业字体(可以使你的论文更加漂亮！)，这些中文字体（宋，黑，仿宋，楷体）可以从
% \href{http://dl.getdropbox.com/u/857066/adobe_chinese_otf.7z}{这里下载}，
% 如果上述链接不能使用，请搜索\textsc{Adobe Opentype 中文字体}自行下载。
% 英文字体使用Windows自带。起始更推荐几款Times（Arial）类似的OTF英文字体，可以使用
% 更多排版、段落的字体特性。
% \item[粗宋] 模板中在需要宋体加黑的地方需要使用\textbf{华文中宋}, 即STZHONGS.TTF。
% \item[xeCJK] 无网络环境中，C\TeX{}完整版和\TeX{}live最新版都包括了需要的xeCJK版本。
% \end{description}
%
% \subsection{使用模板}
% \label{sec:install-cls}
% {\hei 注：默认的发行版本已经包含了可以使用的模板环境，
% 包括编译好的cls以及论文样例源文件，
% 想快速上手的话，可以直接参看\verb|thesis.tex|，进行修改。
% 写作的过程就是将你的论文的内容放到\verb|data|文件夹中，
% 图片放到\verb|figures|文件夹中，用\textsc{jabref}修改\verb|thesis.bib|即可。}
%
% 当用户需要编译生成自己的PDF版论文时，需要依次输入：（注意了，如果不是使用nomencl，
% 则无需使用第二个命令）
% \begin{shell}
% $ xelatex thesis
% $ # makeindex -s nomencl.ist -o thesis.nls thesis.nlo
% $ bibtex thesis
% $ bibtex thesis
% $ xelatex thesis
% $ xelatex thesis
% \end{shell}
%
% 而为了简化用户使用，模板中提供了快捷脚本文件：
% \begin{shell}
% # 下面命令可以直接生成thesis.pdf，你可能只需要这步
% C:\> makepdf.bat
% # linux用户可以直接使用makefile
% $ make pdf
% \end{shell}
% 现在，就要进入激动人心的写作过程了。
%
% \section{使用说明}
% \label{sec:how-to-use}
% 首先，一篇论文（电子信息工程专业为例），主要的构成就是
% 封面导言，正文，表格，图片，公式，
% 交叉引用及文献索引这五部分，下面将分别详细讲解。
%
% \label{sec:howtoask}
% 在开始之前，先问自己几个问题：
% \begin{compactenum}
% \item 我是不是已经掌握了 \LaTeX{} 基础知识？
% \item 我是不是认真地阅读了模板文档？
% \item 周围有没有同学可以帮我？
% \end{compactenum}
% 更推荐用户去阅读示例文档的源代码，改写会给你一个快速的开始。
%
% \subsection{示例文件}
% \label{sec:example}
% 该示例文件是顶层的文件，包括论文属性设置、章节的安排、参考文献附录等。
% 细节用户可以参考\verb|thesis.tex|和\verb|data/|文件夹。
%
% \subsubsection{模板选项}
%
% 论文的第一句话是调用模板：
% \changes{v2.0}{2010/11/10}{增加盲评的说明}
%
%    \begin{macrocode}
%<thesis>%1. 规范硕士导言
%<thesis>% \documentclass[master,ttf]{nudtpaper}
%<thesis>%2. 规范博士导言
%<thesis>% \documentclass[doctor,twoside,ttf]{nudtpaper}
%<thesis>%3. 如果使用是Vista
%<thesis>% \documentclass[master,ttf,vista]{nudtpaper}
%<thesis>%4. 建议使用OTF字体获得较好的页面显示效果
%<thesis>%   OTF字体从网上获得，各个系统名称统一，不用加vista选项
%<thesis>%   如果你下载的是最新的(1201)OTF英文字体，建议修改nudtpaper.cls，使用
%<thesis>%   Times New Roman PS Std
%<thesis>% \documentclass[doctor,twoside,otf]{nudtpaper}
%<thesis>%   另外，新版的论文模板提供了方正字体选项FZ，效果也不错哦
%<thesis>% \documentclass[doctor,twoside,fz]{nudtpaper}
%<thesis>%5. 如果想生成盲评，传递anon即可，仍需修改个人成果部分
%<thesis>% \documentclass[master,otf,anon]{nudtpaper}
%<thesis>%
%    \end{macrocode}
%    \begin{macrocode}
%<*thesis>
\documentclass[master,otf]{nudtpaper}
\usepackage{mynudt}

%</thesis>
%    \end{macrocode}
%
% 模板的参数设置(开关)描述见：
%
%\begin{description}
%\item[master,doctor]
% 硕士论文用master，博士论文就用doctor
%\item[twoside]
% 指定论文为单面打印还是双面打印，当使用\verb|twoside|选项之后，
% 论文会将章节开在奇数页右手边，默认为\verb|openany|单面打印。
%\item[ttf,otf]
% 决定使用何种字体，TTF默认使用Windows自带的字体，而OTF则使用Adobe的字体（需要下载），
% TTF字体的优势是满足学校论文对于字体的要求，缺点是制作出来的PDF文件在浏览时可能发虚，
% 而OTF字体屏幕显示饱满，而且字体有很多选项可以方便\XeTeX{}排版。推荐使用\textbf{otf}
% 选项。不论何种选项，都需要安装宋体中宋(STZHONGSONG)字体(Windows自带)。
%\item[vista]
% 使用\textsc{vista}的用户当调用模板的TTF字体时，系统默认的楷体、仿宋名称是KaiTi和FangSong，
% 而不是KaiTi\_GB2312，这里加入开关进行切换。
%\item[anon]
% 是否为盲评版本，如需盲评，请加上anon。
%\end{description}
%
% 如果需要使用自己定义的命令、宏包，请放于\verb|mynudt.sty|中。
% 事实上，该文件中已经添加了很多有用的宏包和命令，你可以参照修改。
% 这些之所以没有放到模板中，一则为了简洁，二则赋予用户在格式之外更多的自由。
% 里面的宏包有：代码高亮、算法环境、向量命令等，请仔细查看。
%
% 样例文件默认的是硕士论文(master)，OTF字体(otf)。
%
% \subsubsection{封面导言}
% 官方模板中设计论文题目、作者等信息可以跟填空一样完成：
%
% \begin{description}
% \item[论文封头]
% 主要有四部分内容，中图分类号，学号，论文密级和UDC。
% 密级分为：\textbf{秘密} 或者 \textbf{公开}。
%    \begin{macrocode}
%<*thesis>
\classification{TP957}
\serialno{0123456}
\confidentiality{公开}
\UDC{}
%</thesis>
%    \end{macrocode}
%
% \item[论文题目，作者，日期]
% 分别包括中文和英文两部分，由于论文题目可能超过1行，
% 我们提供额外的一个命令\verb|\displaytitle|用来在
% 授权书中填入(限定为)单行的题目; 中文日期需要中文输入大写，英文日期为月年，
% 在论文最终完成后，请\textbf{手动}设定日期。
%
%    \begin{macrocode}
%<*thesis>
\title{国防科大学位论文\LaTeX{}模板\\
使用手册}
\displaytitle{国防科学技术大学学位论文\LaTeX{}模板}
\author{张三}
\zhdate{\zhtoday}
\entitle{How to Use the \LaTeX{} Document Class for NUDT Dissertations}
\enauthor{Zhang San}
\endate{\entoday}
%</thesis>
%    \end{macrocode}
%
% \item[论文分类及其他]
% 主要是作者的学科类别，研究方向，导师信息等。
% 每一项都包括中英文信息：
%
%    \begin{macrocode}
%<*thesis>
\subject{通信与信息工程}
\ensubject{Information and Communication Engineering}
\researchfield{自动目标识别与模糊工程}
\supervisor{李四\quad{}教授}
\cosupervisor{王五\quad{}副教授} % 没有就空着
\ensupervisor{Professor Li Si}
\encosupervisor{}
\papertype{工学}
\enpapertype{Engineering}
%</thesis>
%    \end{macrocode}
%
% \item[中英文摘要]
% 论文中需要写中文以及英文摘要，页码为小写罗马字母，关键字为黑体，
% 英文关键字为\verb|Arial|，
% 模板中定义了相关环境\verb|\cabstract|以及\verb|\eabstract|来书写摘要，
% 以及\verb|\ckeywords|以及\verb|\ekeywords|来写关键字。
% 建议用户将摘要单独放在在\verb|abstract.tex|文件中，
% 在正文中\verb|\begin{cabstract}
分层存储是计算机存储领域的一项重要技术，其核心设计是将数据存储在多层级的存储介质中，通过热点文件识别和数据迁移技术来掩盖访问延迟以及增加吞吐率。分层存储管理的本质是准确、实时的文件分类，当存储层次较多时，可以转化为相邻层级之间热点数据和冷数据的二分类任务。文件分类的准确率十分依赖于对应用I/O行为的理解。本文工作主要包括以下两部分：

%1. 文件静态关联分析。我们提出了两条假设。第一条假设是POSIX文件系统的目录结构具有与自然语言相似的性质，即分布式语义假设。基于该假设，我们使用词嵌入模型将POSIX文件系统中的文件名、目录名映射为高维向量，以量化的方式分析它们之间的语义关联。第二条假设是一条完整路径的各级子向量可以线性相加，得到的结果可以近似表示该文件在文件系统空间中的位置。基于该假设我们提出了路径向量的概念。
文件关联分析。数据块关联性广泛存在于存储系统中，以往的研究通常从POSIX文件目录的树形结构入手，无法充分挖掘语义上的关联性。本文使用词嵌入模型将POSIX文件系统中的文件和目录映射为高维向量，将文件或目录之间的语义关联转化为向量空间中的内积计算，并通过降维可视化的方式予以直观呈现。
%2. I/O行为分析。在上述静态分析的基础上，我们提出了第三条假设：应用的I/O行为研究，可以转化为路径向量的序列分析问题。基于该假设，我们建立了循环神经网络来分析路径向量序列，最终用于解决文件的冷热分类任务，并利用GlusterFS中的tiring模块设计了一套智能分层存储管理的初步方案。

基于I/O行为分析的冷热文件分类。目前，该领域的研究通常局限于短期I/O请求序列的分析和预测，难以在较长时间跨度上挖掘应用的I/O行为模式。本文在文件和目录向量化表示的基础上，使用循环神经网络分析I/O请求序列，最终用于解决文件的冷热分类任务，并利GlusterFS文件系统中的tiring模块设计了一套智能分层存储管理的初步方案。

为验证模型的合理性，本文以开源工程代码编译为工作负载设计了文件冷热分类实验。实验结果表明，对于单进程编译任务，在经过参数调优后可以达到较高的热点文件识别率，且冷数据被误判的概率控制在较低水平。

\end{cabstract}
\ckeywords{分层存储；数据迁移；访问模式；词嵌入；循环神经网络}

\begin{eabstract}
Tierd storage is an import technology in the field of computer storage, which stores data in multiple tiers, masking latency and increasing throughput through file classification and data migration techniques. In fact, the main purpose of hierarchical storage management is to correctly classify files in runtime. When dealing with multi-tired storage system, this mission can be splitted into several sub-missions of binary classification. However, the correctness of file classification heavily depends on the understanding of the I/O behavior. Our research focuses on two aspects.

%Analysis on file correlations. In this part, we propose two assumptions. The first is the distributional hypothesis that the directory structure of POSIX file system is similar to natural language. Based on this hypothesis, we utilize word embedding model to map file names and directory names to high-dimension vectors, analyzing the semantic correlations among those names in a quantified manner. The other is that we could sum up those sub-vectors of a full path to approximately represent this path in the space of file system. Thus, we propose a concept named path vector. 

Analysis on file correlations. File correlations are common in storage systems. Recent researches in this area often focus on the tree structured directory of POSIX file systems which were unable to reveal the implicit semantic correlations among files and directories. We utilize word embedding model to map file names and directory names to high-dimension vectors, analyzing the semantic correlations among them by dot products. With the word embedding model trained, we show the semantic correlations with PCA and visualization.

%Sequential analysis on I/O behavior. Based on aforementioned works, we propose the third assumption, that the research on I/O behavior could be reframed to sequential analysis on path vectors. We build an RNN model to achieve this goal and design a hierarchical storage management solution based on tiering translator in GlusterFS.

File classification based on sequential analysis on I/O behavior. Recent works were not capable of long-term analysis. In this paper, based on aforementioned works, we build an RNN model to analyze the I/O behavior and design a hierarchical storage management solution based on tiering translator in GlusterFS.

To verify these models, we designed file calssification experiments with compiling workloads. The results show that our model performs well in single-process compiling, that the accuracy of hotspot indentification(recall rate) is good with fine-tuned hyperparameters, and the false positive rate controlled bellow.
\end{eabstract}
\ekeywords{Tiered Storage, Data Migration, Access Pattern, Word Embedding, RNN}

|即可。其格式为：
%
% \begin{example}
% \begin{cabstract}
% 中文摘要
% \end{cabstract}
% \ckeywords{关键字}
%
% \begin{eabstract}
% Abstract
% \end{eabstract}
% \ekeywords{Key}
% \end{example}
% \end{description}
%
% \subsubsection{框架构成}
%
% 在定义完论文元素之后，就可以开始写论文正文了。用\LaTeX{}写论文的文件目录构成
% 可以很随意，模板中将图形文件单独放到一个目录中\verb|figure|中，论文正文各个
% 章节置于\verb|data|中；当然也以以\verb|chapter|为目录。
%\changes{v2.2}{2011/05/27}{使用nomencl包管理符号列表}
%\changes{v2.2}{2011/07/08}{默认使用nomencl管理参考文献}
%\changes{v2.2}{2011/09/10}{回复原先使用的denotation方式添加符号列表}
% 
% 如果使用nomencl制作符号列表，在文档开始前要加入\verb|\makenomenclature|命令
% 默认还是使用denotation的方式。nomencl可以参考第二章相关章节，而denote方式
% 请参考\verb|data/denotation.tex|文件（简单的列表环境）。
% 
%<thesis>% 加入makenomenclature命令可用nomencl制作符号列表。
%
%    \begin{macrocode}
%<*thesis> 

\begin{document}
\graphicspath{{figures/}}
%</thesis>
%    \end{macrocode}
% 制作完封面后就是正文四大部分了，分别为：
%
% \begin{compactenum}
% \item frontmatter: 生成目录，图目录，表目录
% \item midmatter: 摘要，符号列表
% \item mainmatter: 正文，致谢，文献，成果
% \item backmatter: 附录
% \end{compactenum}
%
%<thesis>% 制作封面，生成目录，插入摘要，插入符号列表 \\
%<thesis>% 默认符号列表使用denotation.tex，如果要使用nomencl \\
%<thesis>% 需要注释掉denotation，并取消下面两个命令的注释。 \\
%<thesis>% cleardoublepage% \\
%<thesis>% printnomenclature% \\
%
%    \begin{macrocode}
%<*thesis>
\maketitle
\frontmatter
\tableofcontents
\listoftables
\listoffigures

\midmatter
\begin{cabstract}
分层存储是计算机存储领域的一项重要技术，其核心设计是将数据存储在多层级的存储介质中，通过热点文件识别和数据迁移技术来掩盖访问延迟以及增加吞吐率。分层存储管理的本质是准确、实时的文件分类，当存储层次较多时，可以转化为相邻层级之间热点数据和冷数据的二分类任务。文件分类的准确率十分依赖于对应用I/O行为的理解。本文工作主要包括以下两部分：

%1. 文件静态关联分析。我们提出了两条假设。第一条假设是POSIX文件系统的目录结构具有与自然语言相似的性质，即分布式语义假设。基于该假设，我们使用词嵌入模型将POSIX文件系统中的文件名、目录名映射为高维向量，以量化的方式分析它们之间的语义关联。第二条假设是一条完整路径的各级子向量可以线性相加，得到的结果可以近似表示该文件在文件系统空间中的位置。基于该假设我们提出了路径向量的概念。
文件关联分析。数据块关联性广泛存在于存储系统中，以往的研究通常从POSIX文件目录的树形结构入手，无法充分挖掘语义上的关联性。本文使用词嵌入模型将POSIX文件系统中的文件和目录映射为高维向量，将文件或目录之间的语义关联转化为向量空间中的内积计算，并通过降维可视化的方式予以直观呈现。
%2. I/O行为分析。在上述静态分析的基础上，我们提出了第三条假设：应用的I/O行为研究，可以转化为路径向量的序列分析问题。基于该假设，我们建立了循环神经网络来分析路径向量序列，最终用于解决文件的冷热分类任务，并利用GlusterFS中的tiring模块设计了一套智能分层存储管理的初步方案。

基于I/O行为分析的冷热文件分类。目前，该领域的研究通常局限于短期I/O请求序列的分析和预测，难以在较长时间跨度上挖掘应用的I/O行为模式。本文在文件和目录向量化表示的基础上，使用循环神经网络分析I/O请求序列，最终用于解决文件的冷热分类任务，并利GlusterFS文件系统中的tiring模块设计了一套智能分层存储管理的初步方案。

为验证模型的合理性，本文以开源工程代码编译为工作负载设计了文件冷热分类实验。实验结果表明，对于单进程编译任务，在经过参数调优后可以达到较高的热点文件识别率，且冷数据被误判的概率控制在较低水平。

\end{cabstract}
\ckeywords{分层存储；数据迁移；访问模式；词嵌入；循环神经网络}

\begin{eabstract}
Tierd storage is an import technology in the field of computer storage, which stores data in multiple tiers, masking latency and increasing throughput through file classification and data migration techniques. In fact, the main purpose of hierarchical storage management is to correctly classify files in runtime. When dealing with multi-tired storage system, this mission can be splitted into several sub-missions of binary classification. However, the correctness of file classification heavily depends on the understanding of the I/O behavior. Our research focuses on two aspects.

%Analysis on file correlations. In this part, we propose two assumptions. The first is the distributional hypothesis that the directory structure of POSIX file system is similar to natural language. Based on this hypothesis, we utilize word embedding model to map file names and directory names to high-dimension vectors, analyzing the semantic correlations among those names in a quantified manner. The other is that we could sum up those sub-vectors of a full path to approximately represent this path in the space of file system. Thus, we propose a concept named path vector. 

Analysis on file correlations. File correlations are common in storage systems. Recent researches in this area often focus on the tree structured directory of POSIX file systems which were unable to reveal the implicit semantic correlations among files and directories. We utilize word embedding model to map file names and directory names to high-dimension vectors, analyzing the semantic correlations among them by dot products. With the word embedding model trained, we show the semantic correlations with PCA and visualization.

%Sequential analysis on I/O behavior. Based on aforementioned works, we propose the third assumption, that the research on I/O behavior could be reframed to sequential analysis on path vectors. We build an RNN model to achieve this goal and design a hierarchical storage management solution based on tiering translator in GlusterFS.

File classification based on sequential analysis on I/O behavior. Recent works were not capable of long-term analysis. In this paper, based on aforementioned works, we build an RNN model to analyze the I/O behavior and design a hierarchical storage management solution based on tiering translator in GlusterFS.

To verify these models, we designed file calssification experiments with compiling workloads. The results show that our model performs well in single-process compiling, that the accuracy of hotspot indentification(recall rate) is good with fine-tuned hyperparameters, and the false positive rate controlled bellow.
\end{eabstract}
\ekeywords{Tiered Storage, Data Migration, Access Pattern, Word Embedding, RNN}


\chapter*{符号使用说明}
% 可以根据需要在chapter后加星星/去掉星星

\begin{denotation}

\item[HPC] 高性能计算 (High Performance Computing)
\item[cluster] 集群
\item[Itanium] 安腾
\item[SMP] 对称多处理
\item[API] 应用程序编程接口
\item[PI]	聚酰亚胺
\item[MPI]	聚酰亚胺模型化合物，N-苯基邻苯酰亚胺
\item[PBI]	聚苯并咪唑
\item[MPBI]	聚苯并咪唑模型化合物，N-苯基苯并咪唑
\item[PY]	聚吡咙
\item[PMDA-BDA]	均苯四酸二酐与联苯四胺合成的聚吡咙薄膜
\item[$\Delta G$]  	活化自由能~(Activation Free Energy)
\item [$\chi$] 传输系数~(Transmission Coefficient)
\item[$E$] 能量
\item[$m$] 质量
\item[$c$] 光速
\item[$P$] 概率
\item[$T$] 时间
\item[$v$] 速度

\end{denotation}


%</thesis>
%    \end{macrocode}
%
%<thesis>%书写正文，可以根据需要增添章节; 正文还包括致谢，参考文献与成果
%\changes{v1.4}{2009/10/31}{将成果移动到参考文献之后}
%
%    \begin{macrocode}
%<*thesis>
\mainmatter
\chapter{绪论}

\section{研究背景}


\begin{figure}[htp]
\centering
\includegraphics{picmain}
\caption{图 1.1 名称}
\end{figure}
\subsubsection{（1.1.1.1 题目）}
正文内容

正文内容

正文内容

\subsubsection{（1.1.1.2 题目）}
正文内容

正文内容

正文内容

\subsection{（1.1.2 题目）}
正文内容

正文内容

\begin{figure}[htp]
\centering
\includegraphics{picmain}
\caption{图 1.2 名称}
\end{figure}

\section{（1.2 题目）}
正文内容

正文内容

\begin{table}[htp]
\centering
 \begin{minipage}[t]{0.8\linewidth} % 如果想在表格中使用脚注，minipage是个不错的办法
  \caption[表 1.1 名称]{}
	\begin{tabular*}{\textwidth}{lp{10cm}}
      \toprule[1.5pt]
      {\hei 列1} & {\hei 列2} \\
      \midrule[1pt]
      &  \\
      & \\
      & \\
      & \\
      & \\
      & \\
      \bottomrule[1.5pt]
	\end{tabular*}
  \end{minipage}
\end{table}

正文内容

正文内容

正文内容

正文内容

\section{（1.3 题目）}
正文内容

正文内容

正文内容

正文内容

正文内容

正文内容

\subsection{（1.3.1 题目）}
正文内容

\begin{figure}[htp]
\centering
\includegraphics{picmain}
\caption{图 1.3 名称}
\end{figure}

\subsection{（1.3.2 题目）}
正文内容

正文内容

\begin{table}[htp]
\centering
 \begin{minipage}[t]{0.8\linewidth} % 如果想在表格中使用脚注，minipage是个不错的办法
  \caption[表 1.2 名称]{}
	\begin{tabular*}{\textwidth}{lp{10cm}}
      \toprule[1.5pt]
      {\hei 列1} & {\hei 列2} \\
      \midrule[1pt]
      &  \\
      & \\
      & \\
      & \\
      & \\
      & \\
      \bottomrule[1.5pt]
	\end{tabular*}
  \end{minipage}
\end{table}


\chapter{论文正文}
\label{chap:main}
本章将进入论文排版的正文, 按元素分主要包括：
{\kai 字体段落，图片表格，公式定理，参考文献}这几部分。
这个样例文件将包括模板中使用到的所有格式、模板中自定义命令到或者特有的东西，
都将被一一介绍，希望大家在排版自己的学位论文前能细致的看一遍，记住样例的格式和
方法，方便上手。

\section{字体段落}
\label{sec:font}

陈赓（1903年2月27日－1961年3月16日），原名陈庶康，中国湖南湘乡人，军事家。出生将门，其祖父为湘军将领陈翼怀。

Adobe中文字体有四种：

{\kai 楷体\verb|\kai|：陈赓，中国湖南湘乡人，军事家。出生将门，其祖父为湘军将领陈翼怀。%
1952年筹办并任人民解放军军事工程学院第一任院长兼政委，培养国防科技人才。1955年被授予大将军衔。}

{\fs 仿宋\verb|\fs|：陈赓，中国湖南湘乡人，军事家。出生将门，其祖父为湘军将领陈翼怀。%
1952年筹办并任人民解放军军事工程学院第一任院长兼政委，培养国防科技人才。1955年被授予大将军衔。}

{\hei 黑体\verb|\hei|：陈赓，中国湖南湘乡人，军事家。出生将门，其祖父为湘军将领陈翼怀。%
1952年筹办并任人民解放军军事工程学院第一任院长兼政委，培养国防科技人才。1955年被授予大将军衔。}

宋体就是正文字体了。下面测试字体大小，\LaTeX{}默认的列表环境会在
条目之间插入过多的行距，在下面这种情况可能正好，若用户需要
{\kai 正文行距}的列表环境，可以使用compactitem环境，记住这点很重要，不要再
用那种自己修改\verb|itemsep|的傻傻的办法了。
\begin{itemize}
\item[初号] {\song\chuhao 陈赓大将}
\item[小初] {\song\xiaochu 陈赓大将}
\item[一号] {\song\yihao 陈赓大将}
\item[小一] {\song\xiaoyi 陈赓大将}
\item[二号] {\song\erhao 陈赓大将}
\item[小二] {\song\xiaoer 陈赓大将}
\item[三号] {\song\sanhao 陈赓大将}
\item[小三] {\song\xiaosan 陈赓大将}
\item[四号] {\song\sihao 陈赓大将}
\item[小四] {\song\xiaosi 陈赓大将}
\item[五号] {\song\wuhao 陈赓大将}
\item[小五] {\song\xiaowu 陈赓大将}
\end{itemize}

\section{表格明细}
\label{sec:figure}
表格是论文的重要组成部分，我们从简单的表格讲起，到复杂的表格为止。

模板中关于表格的宏包有三个： \textsf{booktabs}、\textsf{array} 和
\textsf{longtabular}。三线表建议使用\textsf{booktabs}中提供的，
包含toprule、midrule 和 bottomrule三条命令，简单干脆！
它们与\textsf{longtable} 能很好的配合使用。下面来看一个表格实例：
\begin{table}[htb]
  \centering
  \begin{minipage}[t]{0.8\linewidth} % 如果想在表格中使用脚注，minipage是个不错的办法
  \caption[模板文件]{模板文件。如果表格的标题很长，那么在表格索引中就会很不美
    观，所以要像 chapter 那样在前面用中括号写一个简短的标题。这个标题会出现在索
    引中。}
  \label{tab:template-files}
    \begin{tabular*}{\linewidth}{lp{10cm}}
      \toprule[1.5pt]
      {\hei 文件名} & {\hei 描述} \\
      \midrule[1pt]
      nudtpaper.ins & \LaTeX{} 安装文件，docstrip\footnote{表格中的脚注} \\
      nudtpaper.dtx & 所有的一切都在这里面\footnote{再来一个}。\\
      nudtpaper.cls & 模板类文件。\\
      nudtpaper.cfg & 模板配置文。cls 和 cfg 由前两个文件生成。\\
      bstutf8.bst   & 参考文献 Bibtex 样式文件。\\
      mynudt.sty    & 常用的包和命令写在这里，减轻主文件的负担。\\
      \bottomrule[1.5pt]
    \end{tabular*}
  \end{minipage}
\end{table}

表 \ref{tab:template-files} 列举了本模板主要文件及其功能，基本上来说论文
中最可能用到的就是这种表格形式了。
请大家注意三线表中各条线对应的命令。这个例子还展示了如何在表格中正确使用脚注。
如果你不需要在表格中插入脚注，可以将minipage环境去掉。
由于\LaTeX{}本身不支持在表格中使用\verb|\footnote|，所以我们不得不将表格放在
小页中，而且最好将表格的宽度设置为小页的宽度，这样脚注看起来才更美观。

另外六院的同学在使用模板时需要使用一种固定宽度（往往是页宽，下面的例子由
rongdonghu提供）的表格，内容需要居中且可以自动调整。
解决办法是自定义了一种\verb|tabularx|中的\textbf{Z}环境，在论文模板中，
该命令已添加到\verb|mynudt.sty|中。下面是这种情况的实例：

\begin{table}[htbp]
\centering
\begin{minipage}[t]{0.9\linewidth}
\caption{Reed Solomon码的典型应用}
\label{tab:RSuse}
\begin{tabularx}{\linewidth}{cZ}
\toprule[1.5pt]
{\hei 应用领域} & {\hei 编码方案}\\
\midrule[1pt]
磁盘驱动器 & RS(32,28,5)码 \footnote{码长为32、维数为28、最小距离为5} \\
CD & 交叉交织RS码(CIRC) \\
DVD & RS(208,192,17)码、RS(182,172,11)码 \\
光纤通信 & RS(255,229,17)码 \\
\bottomrule[1.5pt]
\end{tabularx}
\end{minipage}
\end{table}

我们经常会在表格下方标注数据来源，或者对表格里面的条目进行解释。前面的脚注是一种
不错的方法，如果你不喜欢minipage方法的脚注。
那么完全可以在表格后面自己写注释，比如表~\ref{tab:tabexamp1}。
\begin{table}[htbp]
  \centering
  \caption{复杂表格示例 1}
  \label{tab:tabexamp1}
  \begin{minipage}[t]{0.8\textwidth} 
    \begin{tabularx}{\linewidth}{|l|X|X|X|X|}
      \hline
      \multirow{2}*{\backslashbox{x}{y}}  & \multicolumn{2}{c|}{First Half} & \multicolumn{2}{c|}{Second Half}\\
      \cline{2-5}
      & 1st Qtr &2nd Qtr&3rd Qtr&4th Qtr \\ 
      \hline
      \multirow{2}*{East$^{*}$} &   20.4&   27.4&   90&     20.4 \\
       &   30.6 &   38.6 &   34.6 &  31.6 \\ 
      West$^{**}$ &   30.6 &   38.6 &   34.6 &  31.6 \\ 
      \hline
    \end{tabularx}\\[2pt]
    \footnotesize
    *：东部\\
    **：西部
  \end{minipage}
\end{table}

此外，表~\ref{tab:tabexamp1} 同时还演示了另外三个功能：1）通过 \textsf{tabularx} 的
 \texttt{|X|} 扩展实现表格内容自动调整；2）通过命令 \verb|\backslashbox| 在表头部分
插入反斜线（WORD中很简单，但\LaTeX{}做表格需要一定的（极大的）想象力）；3）就是
使用\verb|multirow|和\verb|multicolumn|命令。

不可否认 \LaTeX{} 的表格功能没有想象中的那么强大，不过只要你足够认真，足够细致，那么
同样可以排出来非常复杂非常漂亮的表格。可是科技论文中那么复杂表格有什么用呢？
上面那个表格就够用啦。

浮动体的并排放置一般有两种情况：1）二者没有关系，为两个独立的浮动体；2）二者隶属
于同一个浮动体。对表格来说并排表格既可以像表~\ref{tab:parallel1}、表~\ref{tab:parallel2} 
使用小页环境，也可以如表~\ref{tab:subtable}使用子表格来做。
图与表同出一源，后面我们将讲解子图(subfloat)的例子。
\begin{table}[htb]
\centering
\noindent\begin{minipage}{0.45\textwidth}
\centering
\caption{第一个并排子表格}
\label{tab:parallel1}
\begin{tabular}{p{2cm}p{2cm}}
\toprule[1.5pt]
111 & 222 \\\midrule[1pt]
222 & 333 \\\bottomrule[1.5pt]
\end{tabular}
\end{minipage}
\begin{minipage}{0.45\textwidth}
\centering
\caption{第二个并排子表格}
\label{tab:parallel2}
\begin{tabular}{p{2cm}p{2cm}}
\toprule[1.5pt]
111 & 222 \\\midrule[1pt]
222 & 333 \\\bottomrule[1.5pt]
\end{tabular}
\end{minipage}
\end{table}
\begin{table}[htbp]
\centering
\caption{并排子表格}
\label{tab:subtable}
\subfloat[第一个子表格]{
\begin{tabular}{p{2cm}p{2cm}}
\toprule[1.5pt]
111 & 222 \\\midrule[1pt]
222 & 333 \\\bottomrule[1.5pt]
\end{tabular}}\hskip2cm
\subfloat[第二个子表格]{
\begin{tabular}{p{2cm}p{2cm}}
\toprule[1.5pt]
111 & 222 \\\midrule[1pt]
222 & 333 \\\bottomrule[1.5pt]
\end{tabular}}
\end{table}

如果您要排版的表格长度超过一页，那么推荐使用\textsf{longtable}命令。
这里随便敲入一些无关的文字，使得正文看上去不是那么的少。
表~\ref{tab:performance} 就是 \textsf{longtable} 的简单示例。
\begin{longtable}[c]{c*{6}{r}}
\caption{实验数据}\label{tab:performance}\\
\toprule[1.5pt]
 测试程序 & \multicolumn{1}{c}{正常运行} & \multicolumn{1}{c}{同步}
& \multicolumn{1}{c}{检查点}   & \multicolumn{1}{c}{卷回恢复}
& \multicolumn{1}{c}{进程迁移} & \multicolumn{1}{c}{检查点} 	\\
& \multicolumn{1}{c}{时间 (s)} & \multicolumn{1}{c}{时间 (s)}
& \multicolumn{1}{c}{时间 (s)} & \multicolumn{1}{c}{时间 (s)}
& \multicolumn{1}{c}{时间 (s)} &  文件（KB）			\\
\midrule[1pt]%
\endfirsthead%

\multicolumn{7}{c}{续表~\thetable\hskip1em 实验数据}\\

\toprule[1.5pt]
 测试程序 & \multicolumn{1}{c}{正常运行} & \multicolumn{1}{c}{同步} 
& \multicolumn{1}{c}{检查点}   & \multicolumn{1}{c}{卷回恢复}
& \multicolumn{1}{c}{进程迁移} & \multicolumn{1}{c}{检查点} 	\\
& \multicolumn{1}{c}{时间 (s)} & \multicolumn{1}{c}{时间 (s)}
& \multicolumn{1}{c}{时间 (s)} & \multicolumn{1}{c}{时间 (s)}
& \multicolumn{1}{c}{时间 (s)} &  文件（KB）			\\
\midrule[1pt]%
\endhead%
\hline%

\multicolumn{7}{r}{续下页}%

\endfoot%
\endlastfoot%
CG.A.2 & 23.05   & 0.002 & 0.116 & 0.035 & 0.589 & 32491  \\
CG.A.4 & 15.06   & 0.003 & 0.067 & 0.021 & 0.351 & 18211  \\
CG.A.8 & 13.38   & 0.004 & 0.072 & 0.023 & 0.210 & 9890   \\
CG.B.2 & 867.45  & 0.002 & 0.864 & 0.232 & 3.256 & 228562 \\
CG.B.4 & 501.61  & 0.003 & 0.438 & 0.136 & 2.075 & 123862 \\
CG.B.8 & 384.65  & 0.004 & 0.457 & 0.108 & 1.235 & 63777  \\
MG.A.2 & 112.27  & 0.002 & 0.846 & 0.237 & 3.930 & 236473 \\
MG.A.4 & 59.84   & 0.003 & 0.442 & 0.128 & 2.070 & 123875 \\
MG.A.8 & 31.38   & 0.003 & 0.476 & 0.114 & 1.041 & 60627  \\
MG.B.2 & 526.28  & 0.002 & 0.821 & 0.238 & 4.176 & 236635 \\
MG.B.4 & 280.11  & 0.003 & 0.432 & 0.130 & 1.706 & 123793 \\
MG.B.8 & 148.29  & 0.003 & 0.442 & 0.116 & 0.893 & 60600  \\
LU.A.2 & 2116.54 & 0.002 & 0.110 & 0.030 & 0.532 & 28754  \\
LU.A.4 & 1102.50 & 0.002 & 0.069 & 0.017 & 0.255 & 14915  \\
LU.A.8 & 574.47  & 0.003 & 0.067 & 0.016 & 0.192 & 8655   \\
LU.B.2 & 9712.87 & 0.002 & 0.357 & 0.104 & 1.734 & 101975 \\
LU.B.4 & 4757.80 & 0.003 & 0.190 & 0.056 & 0.808 & 53522  \\
LU.B.8 & 2444.05 & 0.004 & 0.222 & 0.057 & 0.548 & 30134  \\
EP.A.2 & 123.81  & 0.002 & 0.010 & 0.003 & 0.074 & 1834   \\
EP.A.4 & 61.92   & 0.003 & 0.011 & 0.004 & 0.073 & 1743   \\
EP.A.8 & 31.06   & 0.004 & 0.017 & 0.005 & 0.073 & 1661   \\
EP.B.2 & 495.49  & 0.001 & 0.009 & 0.003 & 0.196 & 2011   \\
EP.B.4 & 247.69  & 0.002 & 0.012 & 0.004 & 0.122 & 1663   \\
EP.B.8 & 126.74  & 0.003 & 0.017 & 0.005 & 0.083 & 1656   \\
\bottomrule[1.5pt]
\end{longtable}

另外，有的同学不想让某个表格或者图片出现在索引里面，那么请使用命令 \verb|\caption*{}|，
这个命令不会给表格编号，也就是出来的只有标题文字而没有“表~XX”，“图~XX”，否则
索引里面序号{\kai 不连续}就显得不伦不类，这也是 \LaTeX{} 里星号命令默认的规则。

\section{绘图插图}

本模板不再预先装载任何绘图包（如 \textsf{pstricks，pgf} 等），完全由你自己来决定。
个人觉得 \textsf{pgf} 不错，不依赖于 Postscript。此外还有很多针对 \LaTeX{} 的
 GUI 作图工具，如 XFig(jFig), WinFig, Tpx, Ipe, Dia, Inkscape, LaTeXPiX,
jPicEdt 等等。本人强烈推荐\textsf{Ipe}。

一般图形都是处在浮动环境中。之所以称为浮动是指最终排版效果图形的位置不一定与源文
件中的位置对应，这也是刚使用 \LaTeX{} 同学可能遇到的问题。
如果要强制固定浮动图形的位置，请使用 \textsf{float} 宏包，
它提供了 \texttt{[H]}（意思是图片就给我放在这里\textcolor{red}{H}ere）参数，
但是除非特别需要，不建议使用\texttt{[H]}，而是推荐使用\texttt{[htbp]}，
给\LaTeX{}更多选择。比如图~\ref{fig:ipe}。
\begin{figure}[htbp] % use float package if you want it here
  \centering
  \includegraphics[width=3in]{hello}
  \caption{利用IPE制图}
  \label{fig:ipe}
\end{figure}

若子图共用一个计数器，
那么请看图~\ref{fig:big1}，它包含两个小图，分别是图~\ref{fig:subfig1} 
和图~\ref{fig:subfig2}。这里推荐使用\verb|\subfloat|，{\bf 不要再用}
\verb|\subfigure|和\verb|\subtable|。
\begin{figure}[htb]
  \centering%
  \subfloat[第一个小图形]{%
    \label{fig:subfig1}
    \includegraphics[height=2cm]{xh}}\hspace{4em}%
  \subfloat[第二个小图形。如果标题很长的话，它会自动换行，这个 caption 就是这样的例子]{%
    \label{fig:subfig2}
    \includegraphics[height=2cm]{xhh}}
  \caption{包含子图形的大图形}
  \label{fig:big1}
\end{figure}

而下面这个例子显示并排$3\times2$的图片，见图\ref{fig:subfig:3x2}:
\begin{figure}[htb]
\centering
\subfloat[]{\includegraphics[width=.27\textwidth]{typography}} \qquad
\subfloat[]{\includegraphics[width=.27\textwidth]{typography}} \qquad
\subfloat[]{\includegraphics[width=.27\textwidth]{typography}} \qquad
\subfloat[]{\includegraphics[width=.27\textwidth]{typography}} \qquad
\subfloat[]{\includegraphics[width=.27\textwidth]{typography}} \qquad
\subfloat[]{\includegraphics[width=.27\textwidth]{typography}}
\caption{并排图片}
\label{fig:subfig:3x2}
\end{figure}

要注意，图\ref{fig:subfig:3x2}例中
\texttt{qquad}相当于\verb|\hspace{2em}|，也就是2个字符的宽度，约0.08倍页宽，
图片宽度设定为0.27倍页宽是合适的；在该环境中，尽量不要手动换行，所以，不妨自己计算一下！

如果要把编号的两个图形并排，那么小页(minipage)就非常有用了，可以分别参考
图\ref{fig:parallel1}和图\ref{fig:parallel2}。其实这个例子和表格一节中并排
放置的表格一摸一样。
\begin{figure}[htb]
\begin{minipage}{0.48\textwidth}
  \centering
  \includegraphics[height=1.2cm]{xhh}
  \caption{并排第一个图}
  \label{fig:parallel1}
\end{minipage}\hfill
\begin{minipage}{0.48\textwidth}
  \centering
  \includegraphics[height=1.2cm]{xhh}
  \caption{并排第二个图}
  \label{fig:parallel2}
\end{minipage}
\end{figure}

图形就说这么多，因为大家在写论文是遇到的最大问题不是怎么把图插进去，
而是怎样做出专业的、诡异的、震撼的图片来，记得在这时参考前面推荐的那
些工具吧，当然必不可少的是Matlab了，至于如何加入中文标注、支持中文等等
可以上网去查，但这里{\kai 推荐一点}，用好export命令，使得插入图片时尽可能的不要
缩放，保证图文的一致性。

\section{公式定理}
\label{sec:equation}
贝叶斯公式如式~(\ref{equ:chap1:bayes})，其中$p(y|\mathbf{x})$为后验；
$p(\mathbf{x})$为先验；分母$p(\mathbf{x})$ 为归一化因子，这是
实际应用中十分恐怖的一个积分式。
\begin{equation}
\label{equ:chap1:bayes}
p(y|\mathbf{x}) = \frac{p(\mathbf{x},y)}{p(\mathbf{x})}=
\frac{p(\mathbf{x}|y)p(y)}{p(\mathbf{x})} 
\end{equation}

论文里面公式越多，\TeX{} 就越 happy。再看一个 \textsf{amsmath} 的例子：
\newcommand{\envert}[1]{\left\lvert#1\right\rvert} 
\begin{equation}\label{detK2}
\det\mathbf{K}(t=1,t_1,\dots,t_n)=\sum_{I\in\mathbf{n}}(-1)^{\envert{I}}
\prod_{i\in I}t_i\prod_{j\in I}(D_j+\lambda_jt_j)\det\mathbf{A}
^{(\lambda)}(\overline{I}|\overline{I})=0.
\end{equation} 

大家在写公式的时候一定要好好看\textsf{amsmath}的文档，并参考模板中的用法：
\begin{multline*}%\tag{[b]} % 这个出现在索引中的
\int_a^b\biggl\{\int_a^b[f(x)^2g(y)^2+f(y)^2g(x)^2]
 -2f(x)g(x)f(y)g(y)\,dx\biggr\}\,dy \\
 =\int_a^b\biggl\{g(y)^2\int_a^bf^2+f(y)^2
  \int_a^b g^2-2f(y)g(y)\int_a^b fg\biggr\}\,dy
\end{multline*}

再看\ref{equ:split}:
\begin{equation}\label{equ:split}
\begin{split}
C(z) &= [z^n] \biggl[\frac{e^{3/4}}{\sqrt{1-z}} +
e^{-3/4}(1-z)^{1/2} + \frac{e^{-3/4}}{4}(1-z)^{3/2}
+ O\Bigl( (1-z)^{5/2}\Bigr)\biggr] \\
&= \frac{e^{-3/4}}{\sqrt{\pi n}} - \frac{5e^{-3/4}}{8\sqrt{\pi
n^3}} + \frac{e^{-3/4}}{128 \sqrt{\pi n^5}} +
O\biggl(\frac{1}{\sqrt{\pi
n^7}}\biggr)
\end{split}
\end{equation}

当然了，数学中必不可少的是定理和证明：
\begin{theorem}
  \label{chapTSthm:rayleigh solution}
  假定 $X$ 的二阶矩存在:
  \begin{equation}
         O_R(\mathbf{x},F)=\sqrt{\frac{\mathbf{u}_1^T\mathbf{A}\mathbf{u}_1} {\mathbf{u}_1^T\mathbf{B}\mathbf{u}_1}}=\sqrt{\lambda_1},
  \end{equation}
  其中 $\mathbf{A}$ 等于 $(\mathbf{x}-EX)(\mathbf{x}-EX)^T$，$\mathbf{B}$ 表示协方差阵 $E(X-EX)(X-EX)^T$，$\lambda_1$
$\mathbf{u}_1$是$\lambda_1$对应的特征向量，
\end{theorem}

对于希腊符号使用\verb|mathbf|命令可能有些问题，所以建议对符号
用\verb|bm|加粗，记得用\verb|\up<greek>|切换正体符号，下面看几个例子：
\verb|\gamma|斜体代表变量$\gamma$，\verb|\bm{\upgamma}|正体代表向量$\bm{\upgamma}$,
。\verb|\Gamma|正体代表操作符号$\Gamma$，
\verb|\bm{\Gamma}|正体粗体代表矩阵形式$\bm{\Gamma}$，
\verb|\varGamma|斜体代表变量$\varGamma$。另外对于大小写斜体的加粗可以见$\bm{\gamma}$和$\bm{\varGamma}$，
但是这两种科技论文中很少出现，这里只做测试。
非符号普通向量就用\verb|\mathbf|吧：$\mathbf{x}_k,\mathbf{X}_k$。
完整测试如下$\omega,\bm{\omega},\upomega,\bm{\upomega},\Omega,\bm{\Omega},\varOmega,\bm{\varOmega}$。

\begin{proof}
上述优化问题显然是一个Rayleigh商问题。我们有
  \begin{align}
     O_R(\mathbf{x},F)=\sqrt{\frac{\mathbf{u}_1^T\mathbf{A}\mathbf{u}_1} {\mathbf{u}_1^T\mathbf{B}\mathbf{u}_1}}=\sqrt{\lambda_1},
 \end{align}
 其中 $\lambda_1$ 下列广义特征值问题的最大特征值：
$$
\mathbf{A}\mathbf{z}=\lambda\mathbf{B}\mathbf{z}, \mathbf{z}\neq 0.
$$
 $\mathbf{u}_1$ 是 $\lambda_1$对应的特征向量。结论成立。
\end{proof}

下面来看看算法环境的定义和使用。
我们知道，故障诊断的最终目的，是将故障定位到部件，而由于信号--部件依赖矩阵的存在，因此，实质性的工作是找出由故障部件发出异常信号，
不妨称为源异常信号，而如前所述，源异常信号与异常信号依赖矩阵$\mathbf{S_a}$的全零列是存在一一对应的关系的。因此，我们只要获得了$\mathbf{S_a}$的全零列的相关信息，
也就获得了源异常信号的信息，从而能进一步找到故障源。
通过以上分析，我们构造算法\ref{alg53}，用于实现非回路故障诊断。
\begin{algorithm}[htbp]
  \caption{非回路故障诊断算法}
  \label{alg53}
  \begin{algorithmic}[1]
    \REQUIRE 信号--部件依赖矩阵$\mathbf{A}$，信号依赖矩阵$\mathbf{S}$，信号状态向量$\alpha$
    \ENSURE 部件状态向量$\gamma$
    \STATE $\mathbf{P}\leftarrow\left(<\alpha>\right)$
    \STATE $\mathbf{S_{a}}\leftarrow\mathbf{P^T}\mathbf{S}\mathbf{P}$
    \FOR{$i=1$ to $S_a$的阶数$m$}
    \STATE $s_i\leftarrow s_i$的第$i$个行向量
    \ENDFOR
    \STATE $\beta_a\leftarrow\lnot \left(s_1\lor s_2\lor \cdots\lor s_m\right)^T$
    \STATE $\beta\leftarrow\mathbf{P}\beta_a$
    \STATE $\gamma\leftarrow\mathbf{A}\beta$
  \end{algorithmic}
\end{algorithm}

第一类故障回路推理与非回路故障推理是算法基本相同，稍微不同的是$\beta_a$的计算。因为第一类故障回路中的信号全部可能是源异常信号，因此我们不必计算
$\beta_a=\lnot \left(\left[s_1\lor s_2\lor \cdots\lor s_m\right]^T\right)$，而直接取$\beta_a=\underbrace{\left[\begin{array}{cccc}1&1&\cdots&1\end{array}\right]^T}_m$，将$\beta_a$代入
算法\ref{alg53}，有
\[\beta=\mathbf{P}\beta_a=\mathbf{P}\underbrace{\left[\begin{array}{cccc}1&1&\cdots&1\end{array}\right]^T}_m=\alpha\]
因此一类故障回路的推理算法变得相当简单，例如算法\ref{alg54}
\begin{algorithm}[htbp]
  \caption{第一类故障回路诊断算法}
  \label{alg54}
  \begin{algorithmic}[1]
    \REQUIRE 信号--部件依赖矩阵$\mathbf{A}$，信号状态向量$\alpha$
    \ENSURE 部件状态向量$\gamma$
    \STATE $\gamma\leftarrow\mathbf{A}\alpha$
  \end{algorithmic}
\end{algorithm}

\section{参考文献}
\label{sec:bib}
当然参考文献可以直接写 bibitem，虽然费点功夫，但是好控制，各种格式可以自己随意改
写，在nudtpaper里面，建议使用JabRef编辑和管理文献，再结合\verb|bstutf8.bst|，
对中文的支持非常不错，格式也很规范。

本模板推荐使用 BIB\TeX，样式文件为 bstutf8.bst，符合学校的参考文献格式（如专利
等引用未加详细测试）。看看这个例子，关于书的\upcite{tex, companion}，
还有这些\upcite{Krasnogor2004e, clzs, zjsw}，关于杂志的\upcite{ELIDRISSI94,
  MELLINGER96, SHELL02}，硕士论文\upcite{zhubajie, metamori2004}，博士论文
\upcite{shaheshang, FistSystem01}，标准文件\upcite{IEEE-1363}，会议论文\upcite{DPMG,kocher99}，%
技术报告\upcite{NPB2}。中文参考文献\upcite{cnarticle}\textsf{特别注意}，需要在\verb|bibitem|中
增加\verb|language|域并设为\verb|zh|，英文此项可不填，之后由\verb|bstutf8|统一处理
(具体就是决定一些文献在中英文不同环境下的显示格式，如等、etc)。
若使用\verb|JabRef|，则你可按下面步骤来设置：
选择\textsf{Options}$\rightarrow$\textsf{Set Up General Fields}，
在\verb|General:|后加入\verb|language|就可以了。

有时候不想要上标，那么可以这样 \cite{shaheshang}，这个非常重要。

\section{代码高亮}
有些时候我们需要在论文中引入一段代码，用来衬托正文的内容，或者体现关键思路的实现。
在模板中，统一使用\texttt{listings}宏包，并且设置了基本的内容格式，并建议用户只
使用三个接口，分别控制：编程语言，行号以及边框。简洁达意即可，下面分别举例说明。

首先是设定语言，来一个C的，使用的是默认设置：
\begin{lstlisting}[language=C]
void sort(int arr[], int beg, int end)
{
  if (end > beg + 1)
  {
    int piv = arr[beg], l = beg + 1, r = end;
    while (l < r)
    {
      if (arr[l] <= piv)
        l++;
      else
        swap(&arr[l], &arr[--r]);
    }
    swap(&arr[--l], &arr[beg]);
    sort(arr, beg, l);
    sort(arr, r, end);
  }
}
\end{lstlisting}

当我们需要高亮Java代码，不需要行号，不需要边框时，可以：
\begin{lstlisting}[language=Java,numbers=none,frame=none]
// A program to display the message
// "Hello World!" on standard output

public class HelloWorld {
 
   public static void main(String[] args) {
      System.out.println("Hello World!");
   }
      
}   // end of class HelloWorld
\end{lstlisting}

细心的用户可能发现，行号被放在了正文框之外，事实上这样是比较美观的，
如果有些用户希望在正文框架之内布置所有内容，可以：
\begin{lstlisting}[language=perl,xleftmargin=2em,framexleftmargin=1.5em]
#!/usr/bin/perl
print "Hello, world!\n";
\end{lstlisting}

好了，就这么多，\texttt{listings}宏包的功能很强大也很复杂，如果需要自己定制，
可以查看其手册，耐心阅读总会找到答案。
\textbf{注意:} 当前代码环境中文注释的处理还不是很完善，对于注释请妥善处理。
在本模板中，推荐算法环境或者去掉中文的listings代码环境。
如果需要包含中文注释，不要求代码高亮，
就用\texttt{code}环境，这个环境是Verbatim的定制版，简单有效，
调用的是fancyvbr宏包，用户可在mynudt.sty中修改它的外观等等。
这里我们还可以给代码加上标签。
\begin{code}[label=hello.c]
public class HelloWorld {
   public static void main(String[] args) {
      System.out.println("Hello World!");
   }
}   // 世界，你好！
\end{code}

\section{符号列表}

{\hei 前面的话：}{\kai\color{blue} 
2.2版本后默认使用nomencl环境，如果你还是希望使用传统的\verb|definition.tex|，那么只需注释掉
顶层文件中的nomenclature即可。}

符号列表使用的是\verb|nomencl|包，自己简单定制了下，使用方法分为四步：
\begin{compactenum}
\item 将\verb|\makenomenclature|语句放在正文前，即\verb|\begin{document}|前面；
\item 将\verb|\printnomenclature|放在论文中，我在例子中将符号列表放在了英文摘要的
后面，正文第一章的前面，当然，你可以根据自己的需要或者教研室的规范放置在合理的位置上，
为了页面引用的正确，在这句话前面放上\verb|\cleardoublepage|；
\item 使用\verb|\nomenclature|命令在论文的各个位置上添加符号定义，语法后面会讲到；
\item 编译。编译需要首先运行一遍xelatex，之后运行
\begin{code}
makeindex -s nomencl.ist -o thesis.nls thesis.nlo
\end{code}
\end{compactenum}

你可以把这句编译命令放在\verb|makepdf.bat|中第一个\verb|xelatex thesis|下面。然后
双击\verb|makepdf.bat|就可以了，论文模板中已经为你添加上了，如果你强烈不想使用
nomencl环境，只要把它注释掉（前面加\verb|rem|）就可以。
另外，由于我使用的是VIM来编辑\TeX{}代码，具体到每个编辑器（诸如WinEDT，TeXWorks等）
如何设定该命令的快捷按钮，诸位可以搜索网上的教程。

下面简单说明下\verb|\nomenclature|命令，语法为。这里插入一些随机的文字，希望
对你在阅读帮助中的思维没有什么不良的影响。
\begin{code}
\nomenclature[<prefix>]{<symbol>}{<desc>}{<null>}
\end{code}
\verb|nomencl|模板的默认排序方法可能（大多都）不满足要求，
论文模板里，我们通过设定\verb|<prefix>|来实现符号列表的排序。
它分为两部分，比如如\verb|[Aa]|，第一个字母的含义是：
\begin{compactitem}
\item[`A'] 符号归为拉丁字母
\item[`G'] 希腊字母
\item[`X'] 上标
\item[`Z'] 下标
\end{compactitem}
每个标识后边的字幕\verb|a-z|作为当前符号组内的排列顺序，比如$\beta$就可以写成
\verb|[Gb]|，诸如此类。当然你一定注意到了，这个排序分组的设定只是为了记忆
方便，并不是强制的，因此你可以有自己的方案，比如Z是Greek，
R是Roman什么的，只要统一就好，只需记住，组间排列是按字母顺序排的。

注意符号表分四列，前三列的含义与命令中相同，
最后一列是符号定义时所在的页码。效果看例子，对于下式:
\begin{equation}\label{eq:heatflux}
   \dot{Q} = k \cdot A \cdot \Delta T
\end{equation}%
\nomenclature[Aq]{$\dot{Q}$}{heat flux}{}%
\nomenclature[Ak]{$k$}{overall heat transfer coefficient,式\eqref{eq:ohtc}}{}%
\nomenclature[Aa]{$A$}{area}{}%
\nomenclature[Al]{$L$}{length}{}%
\nomenclature[At]{$T$}{temperature}{}%
\nomenclature[At]{$\Delta T$}{temperature difference}{}%
\nomenclature[Gr]{$\gamma$}{中文测试, 以及一句很长的物理意义，很有可能超过当前栏的宽度，主要目的是看一看会不会出现某些异常情况。}{}%

或者:
\begin{equation}\label{eq:ohtc}
    \frac{1}{k} = \left[\frac{1}{\alpha _{\mathrm{i}}\,r_{\mathrm{i}}} +
    \sum^n_{j=1}\frac{1}{\lambda _j}\,
    \ln \frac{r_{\mathrm{a},j}}{r_{\mathrm{i},j}} +
    \frac{1}{\alpha _{\mathrm{a}}\,
    r_{\mathrm{a}}}\right] \cdot r_{\mathrm{reference}}
\end{equation}%
\nomenclature[Ga]{$\alpha$}{convection heat transfer coefficient}{}%
\nomenclature[Zi]{i}{in}{}%
\nomenclature[Gl]{$\lambda$}{thermal conductivity}{}%
\nomenclature[Za]{a}{out}{}%
\nomenclature[Zn]{$n$}{number of walls}{}%
\nomenclature[Zj]{$j$}{running parameter}{}%

{\hei 注意事项：}{\kai 模板中定制的nomencl格式在mynudt.sty中，默认是三栏的，分别是：
``符号''，``定义''，``首次出现页码''，
注意这里的符号列表都没有单位，如果你需要额外的栏输入单位（呵呵，聪明的读者可能看出来
了，\verb|nomenclature|命令最后一个是空的，就是用来让你赋予她各种意义的）。
此时就需要你有一点点动手能力了（其实只要会修改表格就行），
方法很简单，比如需要添加``国际单位制''这一栏，则
\begin{compactenum}
\item 论文中\verb|\nomenclature|命令的第三个参数就让他代表单位，也可留空；
\item 将\verb|mynudt.sty|中longtable的表头添加``国际单位制''几个字，
你也可以取其他的名字，放在那个{\kai 应该出现的}位置上；
\item 由于增加了5个字，就把前面栏的宽度数字减5，同时设定第三栏宽度为5，
注意这一步需要你自己调整，记得不要让表格超出边界就行。
\end{compactenum}
}

\section{中文习惯}
\label{sec:chinese}

对于itermize过大的行间距，用户可以使用compactitem环境来替代，但是模板中不进行默认替代，
因为只有用户真正发现列表不好看才会找到这里，而且在示例文件中，
陈赓大将那个列表环境如果压缩了行距会很不好看。谢谢ZhangLei的建议！

{\hei 一个重要的提示：}
作者自己的定义命令、包等，不要放在模板里面，请放到\verb|mynudt.sty|
中，这样模板时，只要覆盖\verb|nudtpaper.cls|即可。

中文破折号为一个两个字宽垂直居中的直线，输入法直接得到的破折号是两个断开的小短线
（——），这看起来不舒服。所以模板中定义了一个破折号的命令 \verb|\pozhehao|，请看：

厚德博学，强军兴国\hfill \pozhehao{}国防科大校训




%%% Local Variables:
%%% mode: latex
%%% TeX-master: "../main"
%%% End:

\begin{ack}
时光荏苒，我的硕士生涯已接进尾声。这几年的时光既漫长又短暂，其中充满了酸甜苦辣，更有收获和成长。几年来，感谢陪我一起度过美好时光的每位尊敬的老师和亲爱的同学，正是你们的帮助，我才能克服困难，正是你们的指导，我才能解决疑惑，直到学业的顺利完成。

本人的学位论文是在我的两位恩师周恩强教授和刘杰教授的殷切关怀和耐心指导下进行并完成的，衷心感谢他们对我的淳淳教诲和悉心关怀。从课题的选择、项目的实施，直至论文的最终完成，恩师始终给予我耐心的指导和支持，我取得的每一点成绩都凝聚着恩师的汗水和心血。恩师开阔的视野、严谨的治学态度、精益求精的工作作风，深深地感染和激励着我，在此谨向他们致以衷心的感谢和崇高的敬意。
  
感谢实验室的师弟师妹们与我一道分享他们青春的快乐!在此还要对实验室所有师兄弟姐妹们在平时开展相关工作中的支持和帮助一并表示感谢。无论在炎热的夏天，还是寒冷的冬季，他们不辞劳苦地为我提供无私的帮助，没有他们的帮助就没有这篇论文的顺利完成。
  
感谢国防科技大学给我提供的平台和机会。今后我会更加努力钻研学术，也更加清醒地意识到必须做终生学习型的人。感谢我的家人常年对我的支持和理解!他们是最爱我的人，也是我亏欠最多的人，他们默默的奉献是我多年来的支持和动力。
  
最后，我要向百忙之中参与审阅、评议本论文各位老师、向参与本人论文答辩的各位老师表示由衷的感谢!人生的每个阶段都值得好好珍惜，这段美好岁月，因为有你们的关心和帮助，我很幸福。我会更加勤奋学习、认真研究，我会努力做得更好，我想这也是我能给你们的最好的回报吧。把最美好的祝福献给你们，愿永远健康、快乐!
\end{ack}


%</thesis>
%    \end{macrocode}
%
% 在\LaTeX{}下管理参考文献将极其方便，建议使用Jabref生成条目，
% 用\verb|\cite|(其中\verb|upcite|是上标索引)索引即可。
% \verb|refs.bib|是你的参考文献名。
%    \begin{macrocode}
%<*thesis>
\cleardoublepage
\phantomsection
\addcontentsline{toc}{chapter}{参考文献}
\bibliographystyle{bstutf8}
\bibliography{ref/refs}

\begin{resume}
%该论文作者在学期间取得的阶段性成果（学术论文等）已满足我校硕士学位评阅相关要求。为避免阶段性成果信息对专家评价学位论文本身造成干扰，特将论文作者的阶段性成果信息隐去。
  \section*{发表的学术论文} % 发表的和录用的合在一起

  \begin{enumerate}[label={[\arabic*]},itemsep=0pt,parsep=0pt,labelindent=26pt,labelwidth=*,leftmargin=0pt,itemindent=*,align=left]
   %[label=\textbf{[\arabic*]},itemindent=*, align=left] %老版本缩进对齐
   
  %\addtolength{\itemsep}{-.36\baselineskip}%缩小条目之间的间距，下面类似
  \item Hui Chen, Enqiang Zhou, Jie Liu, Zhicheng Zhang, et al. An RNN Based Mechanism for File Prefetching. 18th Distributed Computing and Applications for Business Engineering and Science.

  \end{enumerate}
\end{resume}
%</thesis>
%    \end{macrocode}
%
%<thesis>% 最后，需要的话还要生成附录，全文随之结束。
%    \begin{macrocode}
%<*thesis>
\appendix
\backmatter
% TeX
\chapter{模板提供的希腊字母命令列表}

大写希腊字母：
\begin{table}[htbp]
\centering
\begin{tabular}{llll}
\toprule
$\Gamma$~\verb|\Gamma| & $\Lambda$~\verb|\Lambda| & $\Sigma$~\verb|\Sigma| & $\Psi$~\verb|\Psi| \\
$\Delta$~\verb|\Delta| & $\Xi$~\verb|\Xi| & $\Upsilon$~\verb|\Upsilon| & $\Omega$~\verb|\Omega| \\
$\Theta$~\verb|\Theta| & $\Pi$~\verb|\Pi| & $\Phi$~\verb|\Phi| & \\
\midrule
$\varGamma$~\verb|\varGamma| & $\varLambda$~\verb|\varLambda| & $\varSigma$~\verb|\varSigma| & $\varPsi$~\verb|\varPsi| \\
$\varDelta$~\verb|\varDelta| & $\varXi$~\verb|\varXi| & $\varUpsilon$~\verb|\varUpsilon| & $\varOmega$~\verb|\varOmega| \\
$\varTheta$~\verb|\varTheta| & $\varPi$~\verb|\varPi| & $\varPhi$~\verb|\varPhi| & \\
\bottomrule
\end{tabular}
\end{table}

小写希腊字母：
\begin{table}[htbp]
\centering
\begin{tabular}{llll}
\toprule
$\alpha$~\verb|\alpha| & $\theta$~\verb|\theta| & $o$~\verb|o| & $\tau$~\verb|\tau| \\ 
$\beta$~\verb|\beta| & $\vartheta$~\verb|\vartheta| & $\pi$~\verb|\pi| & $\upsilon$~\verb|\upsilon| \\ 
$\gamma$~\verb|\gamma| & $\iota$~\verb|\iota| & $\varpi$~\verb|\varpi| & $\phi$~\verb|\phi| \\ 
$\delta$~\verb|\delta| & $\kappa$~\verb|\kappa| & $\rho$~\verb|\rho| & $\varphi$~\verb|\varphi| \\ 
$\epsilon$~\verb|\epsilon| & $\lambda$~\verb|\lambda| & $\varrho$~\verb|\varrho| & $\chi$~\verb|\chi| \\ 
$\varepsilon$~\verb|\varepsilon| & $\mu$~\verb|\mu| & $\sigma$~\verb|\sigma| & $\psi$~\verb|\psi| \\ 
$\zeta$~\verb|\zeta| & $\nu$~\verb|\nu| & $\varsigma$~\verb|\varsigma| & $\omega$~\verb|\omega| \\ 
$\eta$~\verb|\eta| & $\xi$~\verb|\xi| & $\varkappa$~\verb|\varkappa| & $\digamma$~\verb|\digamma| \\ 
\midrule
$\upalpha$~\verb|\upalpha| & $\uptheta$~\verb|\uptheta| & $\mathrm{o}$~\verb|\mathrm{o}| & $\uptau$~\verb|\uptau| \\ 
$\upbeta$~\verb|\upbeta| & $\upvartheta$~\verb|\upvartheta| & $\uppi$~\verb|\uppi| & $\upupsilon$~\verb|\upupsilon| \\ 
$\upgamma$~\verb|\upgamma| & $\upiota$~\verb|\upiota| & $\upvarpi$~\verb|\upvarpi| & $\upphi$~\verb|\upphi| \\ 
$\updelta$~\verb|\updelta| & $\upkappa$~\verb|\upkappa| & $\uprho$~\verb|\uprho| & $\upvarphi$~\verb|\upvarphi| \\ 
$\upepsilon$~\verb|\upepsilon| & $\uplambda$~\verb|\uplambda| & $\upvarrho$~\verb|\upvarrho| & $\upchi$~\verb|\upchi| \\ 
$\upvarepsilon$~\verb|\upvarepsilon| & $\upmu$~\verb|\upmu| & $\upsigma$~\verb|\upsigma| & $\uppsi$~\verb|\uppsi| \\ 
$\upzeta$~\verb|\upzeta| & $\upnu$~\verb|\upnu| & $\upvarsigma$~\verb|\upvarsigma| & $\upomega$~\verb|\upomega| \\ 
$\upeta$~\verb|\upeta| & $\upxi$~\verb|\upxi| & & \\ 
\bottomrule
\end{tabular}
\end{table}

希腊字母属于数学符号类别，请用\verb|\bm|命令加粗，其余向量、矩阵可用\verb|\mathbf|。


\end{document}
%</thesis>
%    \end{macrocode}
%
% 当然还有一些收尾工作，校验审阅自不必说。接下来你需要：修改论文中英文日期，
% 生成盲评，生成明(盲)评A3封面。
%
% {\color{blue}Happy \TeX{}ing! 欢迎提各式各样的意见！}
%
% \newpage\relax%
%
% \StopEventually{\PrintChanges}
% \clearpage
%
% \section{实现细节}
% 我们首先介绍文档模板的基本信息以及宏包和配置，
% 然后依照国防科学技术大学论文模板的书写规范一节一节的介绍实现步骤。
%
% \changes{v1.2}{2009/09/28}{添加了A3封面制作}
%
% \subsection{基本信息}
%    \begin{macrocode}
%<cls>\NeedsTeXFormat{LaTeX2e}[1999/12/01]
%<cls>\ProvidesClass{nudtpaper}
%<cfg>\ProvidesFile{nudtpaper.cfg}
%<cls|cfg>[2011/07/17 v2.2 NUDT paper template]
%    \end{macrocode}
%
% \subsection{宏包配置}
%
%<*cls>
%
%\changes{v0.99}{2009/08/17}{add package options}
% 当前的宏包选项在之前已经介绍了，下面是实现步骤，就是几个\verb|if|。
%\changes{v1.6}{2009/12/01}{添加单独的单双面控制}
%\changes{v2.0}{2010/11/09}{添加盲评控制}
%
%    \begin{macrocode}
\newif\ifismaster\ismastertrue
\newif\ifisttf\isttffalse
\newif\ifisfz\isfzfalse
\newif\ifisotf\isotffalse
\DeclareOption{master}{\ismastertrue}
\DeclareOption{doctor}{\ismasterfalse}
\newif\ifisanon\isanonfalse
\DeclareOption{anon}{\isanontrue}
\newif\ifistwoside\istwosidefalse
\DeclareOption{twoside}{\istwosidetrue}
\DeclareOption{ttf}{\isttftrue}
\DeclareOption{fz}{\isfztrue}
\DeclareOption{otf}{\isotftrue}
\newif\ifisvista\isvistafalse
\DeclareOption{vista}{\isvistatrue}
\DeclareOption*{\PackageWarning{nudtpaper}{Unknown Option '\CurrentOption'}}
\ProcessOptions\relax
%    \end{macrocode}
%
% 首先调用在文档类书写中需要的过程控制语句，在计算一些\verb|length|时要用到
%    \begin{macrocode}
\RequirePackage{ifthen,calc}
%    \end{macrocode}
%
% 接着我们导入文本类，该模板基于标准的书籍模板book，其默认格式为单面打印。
% 博士论文如需双面打印，必须指定\verb|twoside|选项。双开的含义是章节总是
% 起在右手边，左手空白页为完全的空白页，不包含页眉页脚。
%
% \changes{v1.6}{2009/12/01}{修改开关选项}
%
%    \begin{macrocode}
\ifistwoside
  \LoadClass[a4paper,12pt,openright,twoside]{book}
\else
  \LoadClass[a4paper,12pt,openany]{book}
\fi
%    \end{macrocode}
%
% 我们直接用\textsf{geometry}宏包进行页面边距的设定，调用titlesec设定标题以及页眉页脚，
% 用\textsf{titletoc}设定目录格式。需要改动的可以参考这三个宏包的说明文档。
%
%    \begin{macrocode}
\RequirePackage[includeheadfoot]{geometry}
\RequirePackage[center,pagestyles]{titlesec}
\RequirePackage{titletoc}
%    \end{macrocode}
%
% 文档中另外重要的两个部分是表格和图片。
% 首先来看图片：\textsf{graphicx}宏包是必不可少的，
% 并排图形。\textsf{subfigure} 已经不再推荐，用新的 \textsf{subfig}。
% 加入 \verb|config| 选项
% 以便兼容 \textsf{subfigure} 的命令。浮动图形和表格标题样式。\textsf{caption2} 已经不
% 推荐使用，采用新的 \textsf{caption}。它会自动被 \textsf{subfig} 装载进来。所以可以在
% 后面使用 \textbf{captionsetup} 命令，宏包\textsf{float}的作用是可以用H命令，
% 将浮动对象强制放在这里(副作用是版面可能不好):
%
%    \begin{macrocode}
\RequirePackage{graphicx}
\RequirePackage[config]{subfig}
\RequirePackage{float}
%    \end{macrocode}
%
% 再来看表格：我们采用\textsf{longtable}来处理长的表格，还需要\textsf{array}包;
% 标准的论文需要表格为三线表，这里引用\textsf{booktabs}宏包来处理，
% 这样，我们就可以简单的使用\verb|\toprule|，\verb|\midrule|，\verb|bottomrulle|
% 这样的命令;
% 为了在表格中支持跨行，需要引入\textsf{multirow}包，\textsf{tabularx}的作用是为了使用
% 固定宽度的表格，\textsf{slashbox}可以让我们在表格中使用反斜线:
%    \begin{macrocode}
\RequirePackage{array}
\RequirePackage{longtable}
\RequirePackage{booktabs}
\RequirePackage{multirow}
\RequirePackage{tabularx}
\RequirePackage{slashbox}
%    \end{macrocode}
% 表格和图片的例子可以搜索C\TeX{}论坛或者看示例文件。
%
% 引入\textsf{paralist}来达到比较好看的列表环境
%    \begin{macrocode}
\RequirePackage[neverdecrease]{paralist}
%    \end{macrocode}
%
% 文档中还需要一定的色彩控制和字体控制
%    \begin{macrocode}
\RequirePackage{xcolor}
%    \end{macrocode}
%
% 为了排出漂亮的数学公式，\textsf{amsmath}包是必不可少的，
% 需要注意的是，新版本的论文模板仍旧使用\textsf{txfonts}宏包，
% 为了支持希腊正体字母，需要调用\verb|upgreek.sty|，使用方法是\verb|\up<greek>|。
% 注意到这个宏包前面加上了\verb|Symbolsmallscale|选项，这是为了配合
% \verb|txfonts|希腊字体的大小而设定的。如果用户不满意这个宏包的积分号
% 等符号，倾向与使用传统的\LaTeX{}风格的数学符号，那么可以使用
% \textsf{mathptmx}宏包，但要把\verb|upgreek|的选项改为\verb|Symbol|，要不然
% 正体希腊字母要显得小一点哦。
% 而大写斜体希腊字母(变量)可以通过\textsf{amsmath}的\verb|\var<Greek>|得到。
% 当然，对于希腊字母的加粗推荐使用\verb|bm|宏包，一般变量的加粗那就使用
% \verb|\mathbf|吧！
% \changes{v2.0}{2010/11/09}{去掉fontspec，传递no-math到xeCJK，加入bm宏包}
% \changes{v2.2}{2011/07/16}{去掉txfonts宏包，使用lm字体，添加svgreek.sty}
% \changes{v2.2}{2011/07/16}{修改，仍旧使用upgreek, mathptmx, bm组合}
% \changes{v2.2}{2011/09/25}{修改，使用upgreek, txfonts, bm组合}
% \changes{v2.2}{2012/11/28}{给用户提供额外的选项，还是mtpro比较漂亮}
% \changes{v2.3}{2013/12/27}{调整数学字体，加入mtpro使用说明，并且仿照IEEE模板，修改displaypenalty}
%    \begin{macrocode}
\RequirePackage{amsmath,amssymb}
\RequirePackage{txfonts}
\RequirePackage[Symbolsmallscale]{upgreek}
%\RequirePackage{amsmath}
%\RequirePackage[amsbb,eufrak,compatiblegreek,subscriptcorrection,nofontinfo]{mtpro2}
\interdisplaylinepenalty=2500
\RequirePackage{bm}
\RequirePackage[T1]{fontenc}
\RequirePackage[amsmath,thmmarks,hyperref]{ntheorem}
%    \end{macrocode}
% 需要注意的是，如果用户有\verb|mtpro2|包，还是强烈建议使用这个的，因为数学公式
% 在这个包下显得特别的美观。虽然下载和安装不属于这篇使用说明的范畴，但是，上面的注释部分
% 可以给大家如何使用的一个简单的例子。当你安装好\verb|mtpro2|之后，主要取消注释，并且将
% 上面的三个包注释掉即可。
%
% 本文档类直接采用\XeTeX{}引擎，方便了字体配置以及编译，
% 这里需要调用\textsf{XeCJK}宏包，no--math的作用是不改变先前数学宏包设定的数学字体。
% 同时采用\textsf{indentfirst}宏包管理文字的缩进:
% \changes{v1.8}{2010/10/15}{修改了默认的xeCJK的选项，为了兼容旧的xeCJK版本，normalindentfirst选项暂不使用，而是在后面添加indentfirst包}
% \changes{v2.0}{2010/11/10}{传递no-math给xeCJK里面的fontspec宏包}
% \changes{v2.2}{2011/07/03}{移除CJKtextspace, CJKmathspace, CJKnumber选项}
%
%    \begin{macrocode}
\RequirePackage[CJKnumber,CJKchecksingle,no-math]{xeCJK}
\RequirePackage{indentfirst}
%    \end{macrocode}
%
% 另外一个关键部分是文献索引，包括书签以及参考文献的索引，记得\textsf{hyperref}配合
% \XeTeX{}使用时暂不能开启Unicode选项，新的发行版已经移除\textsf{hypernat}包。另外还要注意，你最终的打印版肯定
% 不希望有花花绿绿的链接，对吧？那就把下面那行\verb|hyperref|注释掉就行了。
% \changes{v2.1}{2010/12/29}{移除hypernat包}
% \changes{v2.2}{2011/07/17}{移除hyperref的CJKbookmarks旋向}
% \changes{v2.3}{2013/12/27}{加入色彩版hyperref}
%    \begin{macrocode}
\RequirePackage[numbers,sort&compress,square]{natbib}
\RequirePackage[colorlinks=true,linkcolor=blue,citecolor=red,pdfborder=0 1 1]{hyperref}
%\RequirePackage[pdfborder=0 0 1]{hyperref}
%    \end{macrocode}
%</cls>
%
%\subsection{基础配置}
% 本章主要介绍模板中用到的基本的元素和定义，现在包括两部分: 字体，字号和字体命令
%
%\subsubsection{字体定义}
% 我们首先来处理\TeX{}中最令人棘手的字体问题，
% 在使用\textsf{XeCJK}包之后，配置和选择很容易，
% 预先设定好一些字体命令是为了后面方便的更改文本字体的需要。
% 首先我们开启\TeX{}连字符:
%    \begin{macrocode}
%<*cls>
\defaultfontfeatures{Mapping=tex-text}
%</cls>
%    \end{macrocode}
%
% 之后用\textsc{XeCJK}包提供的命令设定字体，用户可以选择使用TTF还是OTF字体，
% Adobe的OpenType字体在排版上更具备优势，文档显示锐利，推荐使用。
% 另外在这一个新版本中，我们推荐用户也可以使用方正的字体，只要使用\verb|FZ|选项即可。
% 中注释掉相关的字体就可以。方正字体的有点是标点符号的位置无需修正，且字体之间配合很好。
% \verb|setcharclass|的作用是纠正xunicode、xeCJK的一些设定:
%
% \changes{v0.99}{2009/08/17}{add options TTF and OTF}
% \changes{v0.993}{2009/08/25}{加入VISTA用户选项}
% \changes{v1.9}{2010/10/28}{定义一个cusong字体，使用的是中宋}
% \changes{v2.3}{2013/12/27}{用户可以考虑使用方正字体，加入FZ选项}
%
%    \begin{macrocode}
%<*cls>
\xeCJKsetcharclass{"0}{"2E7F}{0}
\xeCJKsetcharclass{"2E80}{"FFFF}{1}
\newcommand\installTTF{%
  \setmainfont{Times New Roman}
  \setsansfont{Arial}
  \setmonofont{Courier New}
  \ifisvista
    \setCJKmainfont[BoldFont={SimHei},ItalicFont={KaiTi}]{SimSun}
    \setCJKmonofont{KaiTi} % Pluto use LiSu Thu use Kaiti, orig is SimSun
    \setCJKfamilyfont{fs}{FangSong}
    \setCJKfamilyfont{kai}{KaiTi}
  \else
    \setCJKmainfont[BoldFont={SimHei},ItalicFont={KaiTi_GB2312}]{SimSun}
    \setCJKmonofont{KaiTi_GB2312} % Pluto use LiSu Thu use Kaiti, orig is SimSun
    \setCJKfamilyfont{fs}{FangSong_GB2312}
    \setCJKfamilyfont{kai}{KaiTi_GB2312}
  \fi
  \setCJKsansfont{SimHei}
  \setCJKfamilyfont{song}{SimSun}
  \setCJKfamilyfont{hei}{SimHei}
  \setCJKfamilyfont{li}{LiSu}
  \setCJKfamilyfont{you}{YouYuan}
  \setCJKfamilyfont{cusong}{STZhongsong}
}
\newcommand\installOTF{%
  \setmainfont{Times New Roman} % could be changed to "Times New Roman PS Std" !!
  \setsansfont{Arial}
  \setmonofont{Courier New}
  \setCJKmainfont[BoldFont={Adobe Heiti Std},ItalicFont={Adobe Kaiti Std}]{Adobe Song Std}
  \setCJKsansfont{Adobe Heiti Std}
  \setCJKmonofont{Adobe Kaiti Std}
  \setCJKfamilyfont{song}{Adobe Song Std}
  \setCJKfamilyfont{hei}{Adobe Heiti Std}
  \setCJKfamilyfont{fs}{Adobe Fangsong Std}
  \setCJKfamilyfont{kai}{Adobe Kaiti Std}
  \setCJKfamilyfont{li}{Adobe Kaiti Std}
  \setCJKfamilyfont{you}{Adobe Kaiti Std}
  \setCJKfamilyfont{cusong}{STZhongsong}
}
\newcommand\installFZ{%
  \setmainfont{Times New Roman} % could be changed to "Times New Roman PS Std" !!
  \setsansfont{Arial}
  \setmonofont{Courier New}
  \setCJKmainfont[BoldFont={FZHei-B01},ItalicFont={FZKai-Z03}]{FZShuSong_GB18030-Z01}
  \setCJKsansfont{FZHei-B01} % Hei
  \setCJKmonofont{FZKai-Z03} % Kai
  \setCJKfamilyfont{song}{FZShuSong_GB18030-Z01}
  \setCJKfamilyfont{hei}{FZHei-B01}
  \setCJKfamilyfont{fs}{FZKai-Z03} % Kai
  \setCJKfamilyfont{kai}{FZKai-Z03} % Kai
  \setCJKfamilyfont{li}{FZKai-Z03} % Kai
  \setCJKfamilyfont{you}{FZKai-Z03} % Kai
  \setCJKfamilyfont{cusong}{FZXiaoBiaoSong-B05} % 小标宋
}
\newcommand{\cusong}{\CJKfamily{cusong}} % 中宋作为加粗宋体
%</cls>
%    \end{macrocode}
%
% \changes{v1.6}{2009/12/01}{替换OTF英文字体为标准Windows自带字体}
% \changes{v2.3}{2013/12/27}{添加FZ字体选项}
% 之后我们根据你的设定决定安装什么字体：
%
%    \begin{macrocode}
%<*cls>
\ifisttf
  \installTTF
\else
  \ifisfz
    \installFZ
  \else
    \installOTF
  \fi
\fi
%</cls>
%    \end{macrocode}
%
% 选定好字体之后，就是设定字体别名，这样我们就可以在文档的其他部分直接使用较短的命令来
% 指定特定的字体了:
%
%    \begin{macrocode}
%<*cls>
\newcommand{\song}{\CJKfamily{song}}    % 宋体
\newcommand{\fs}{\CJKfamily{fs}}        % 仿宋体
\newcommand{\kai}{\CJKfamily{kai}}      % 楷体
\newcommand{\hei}{\CJKfamily{hei}}      % 黑体
\newcommand{\li}{\CJKfamily{li}}        % 隶书
\newcommand{\you}{\CJKfamily{you}}      % 幼圆
\def\songti{\song}
\def\fangsong{\fs}
\def\kaishu{\kai}
\def\heiti{\hei}
\def\lishu{\li}
\def\youyuan{\you}
%</cls>
%    \end{macrocode}
%
% \subsubsection{字号定义}
%下面就是定义字号大小，这一部分我们有两个参考，其一是：
%
% \begin{verbatim}
% 参考科学出版社编写的《著译编辑手册》(1994年)
% 七号      5.25pt       1.845mm
% 六号      7.875pt      2.768mm
% 小五      9pt          3.163mm
% 五号      10.5pt       3.69mm
% 小四      12pt         4.2175mm
% 四号      13.75pt      4.83mm
% 三号      15.75pt      5.53mm
% 二号      21pt         7.38mm
% 一号      27.5pt       9.48mm
% 小初      36pt         12.65mm
% 初号      42pt         14.76mm
%
% 这里的 pt 对应的是 1/72.27 inch，也就是 TeX 中的标准 pt
% \end{verbatim}
%
% 另外一个来自WORD中的设定：
% \begin{verbatim}
% 初号 = 42bp = 14.82mm = 42.1575pt
% 小初 = 36bp = 12.70mm = 36.135 pt
% 一号 = 26bp = 9.17mm = 26.0975pt
% 小一 = 24bp = 8.47mm = 24.09pt
% 二号 = 22bp = 7.76mm = 22.0825pt
% 小二 = 18bp = 6.35mm = 18.0675pt
% 三号 = 16bp = 5.64mm = 16.06pt
% 小三 = 15bp = 5.29mm = 15.05625pt
% 四号 = 14bp = 4.94mm = 14.0525pt
% 小四 = 12bp = 4.23mm = 12.045pt
% 五号 = 10.5bp = 3.70mm = 10.59375pt
% 小五 = 9bp = 3.18mm = 9.03375pt
% 六号 = 7.5bp = 2.56mm
% 小六 = 6.5bp = 2.29mm
% 七号 = 5.5bp = 1.94mm
% 八号 = 5bp = 1.76mm
%
% 1bp = 72.27/72 pt
% \end{verbatim}
%
% 我们采用习惯的字号设定方法(也就是WORD中的设定)，首先编写字体设置命令：
%
%\begin{macro}{\choosefont}
% 我们可以使用 |\choosefont| 来选择字体, 字体设定这些大多是从清华的模板拷过来的。
%
%    \begin{macrocode}
%<*cls>
\newlength\thu@linespace
\newcommand{\thu@choosefont}[2]{%
    \setlength{\thu@linespace}{#2*\real{#1}}%
    \fontsize{#2}{\thu@linespace}\selectfont}
\def\thu@define@fontsize#1#2{%
    \expandafter\newcommand\csname #1\endcsname[1][\baselinestretch]{%
    \thu@choosefont{##1}{#2}}}
%</cls>
%    \end{macrocode}
%\end{macro}
%
%设定具体的字体大小:
%
%    \begin{macrocode}
%<*cls>
\thu@define@fontsize{chuhao}{42bp}
\thu@define@fontsize{xiaochu}{36bp}
\thu@define@fontsize{yihao}{26bp}
\thu@define@fontsize{xiaoyi}{24bp}
\thu@define@fontsize{erhao}{22bp}
\thu@define@fontsize{xiaoer}{18bp}
\thu@define@fontsize{sanhao}{16bp}
\thu@define@fontsize{xiaosan}{15bp}
\thu@define@fontsize{sihao}{14bp}
\thu@define@fontsize{banxiaosi}{13bp}
\thu@define@fontsize{xiaosi}{12bp}
\thu@define@fontsize{dawu}{11bp}
\thu@define@fontsize{wuhao}{10.5bp}
\thu@define@fontsize{xiaowu}{9bp}
\thu@define@fontsize{liuhao}{7.5bp}
\thu@define@fontsize{xiaoliu}{6.5bp}
\thu@define@fontsize{qihao}{5.5bp}
\thu@define@fontsize{bahao}{5bp}
%</cls>
%    \end{macrocode}
%
%\subsubsection{自定命令}
% 有一些常量，测试，自定义的命令等都放在这里，待到论文逐渐完善之后再做定夺，
% 当然用户自己的命令也可以在此添加，事实上如果natbib传递的是superscript，
% \verb|cite|命令默认就成了上标了。这里不加入这个选项，而是单独编写一个命令:
%
%    \begin{macrocode}
%<*cls>
\newcommand{\upcite}[1]{\textsuperscript{\cite{#1}}} % 上标形式引用
\newcommand{\china}{中华人民共和国}
\def\nudtpaper{\textsc{Nudt}\textsc{Paper}}
\newcommand{\pozhehao}{\kern0.3ex\rule[0.8ex]{2em}{0.1ex}\kern0.3ex}
%</cls>
%    \end{macrocode}
%
%\subsubsection{中文元素}
%
% 默认的页面元素的英文名，诸如Contents为目录，Abstract为摘要等，
% 我们首先将他们一一中文化：
% \changes{v0.992}{2009/08/19}{修改图表编号格式}
% \changes{v1.3}{2009/10/14}{修改图目录和表目录}
%
%    \begin{macrocode}
%<*cls>
\renewcommand\contentsname{目\hspace{1em}录}
\renewcommand\listfigurename{图\hspace{1em}目\hspace{1em}录}
\renewcommand\listtablename{表\hspace{1em}目\hspace{1em}录}
\newcommand\listequationname{公式索引}
\newcommand\equationname{公式}
\renewcommand\bibname{参考文献}
\renewcommand\indexname{索引}
\renewcommand\figurename{图}
\renewcommand\tablename{表}
\renewcommand\appendixname{附录}
\def\CJK@today{\CJKdigits{\the\year} 年 \CJKnumber{\the\month} 月}
\newcommand\zhtoday{\CJK@today}
\newcommand\entoday{\today{}}
%</cls>
%    \end{macrocode}
%
% 好，下面就开始按照论文模板要求进行排版！
%
%\subsection{编写要求}
% 学校规定，论文需采用白色纸双面打印。
% 学位论文用A4（$210mm\times{}297mm$）标准大小的白纸，
% 在打字或印刷时，要求纸的四周留足空白边缘，以便装订、复制和读者批注。
% 每一面的上方（天头）和下方（地角）分别留边25mm，左侧（订口）
% 和右侧（切口）分别留边30mm，页眉与页脚分别为23mm。
%
% 实现起来很简单，只要调用\textsf{geometry}的版面控制命令即可，
% 方法为先把word模板转化为PDF，
% 用Adobe的裁剪功能查看页边距，进行微调，直到比对正确为止，设定如下：
%
% \changes{v0.991}{2009/08/18}{modify bottom skip}
% \changes{v1.1}{2009/09/26}{修改footskip容限以及bottom的值，为了容下longtab的''下一页''}
% \changes{v1.4}{2009/10/28}{减小页眉skip 1mm，用word叠印}
% \changes{v1.4}{2009/10/30}{增大页眉sep .5mm，用word叠印}
%
%    \begin{macrocode}
%<*cls>
\geometry{top=21mm,bottom=25.5mm,left=30mm,right=30mm}
\geometry{headheight=9mm,headsep=1mm,footskip=9mm}
%</cls>
%    \end{macrocode}
%
%\subsection{页眉页脚}
%
% 我们采用titlesec进行页面配置。
% 页面中的主要元素有Chapter，Section，Subsection等元素的外观，
% 位置，颜色字体等，页面元素还包括页眉页脚。这种方法配置简便，易管理。
% 国防科大的论文需要在页眉处画两根横线，我们通过下面的命令实现:
%
%\begin{macro}{\setheadrule}
% 这个命令属于更改\textsf{titlesec}中的一个画页眉的命令,稍加调整:
% \changes{v0.991}{2009/08/18}{modify headrull, s.t. all geometry match}
% \changes{v1.9}{2010/10/28}{去掉headsep，修改headrule，在sethead后添加raisebox}
%
%    \begin{macrocode}
%<*cls>
\renewcommand\setheadrule[1]{%
  \ifdim#1=\z@
    \let\makeheadrule\@empty
  \else
    \def\makeheadrule{%
    \makebox[0pt][l]{\rule[.2\baselineskip]{\linewidth}{1.5pt}}%
    \rule{\linewidth}{1.5pt}}%
  \fi}
%</cls>
%    \end{macrocode}
%\end{macro}
%
% 由于Chapter第一页默认是\verb|plain|页面格式，
% 章节的其余部分是在Matter中设定的页面格式，为了简单起见，
% 我们就直接更改\verb|plain|页面设置，
% 要求为5号宋体居中放置，画页眉页脚，页脚为1磅黑线
%
% \changes{v0.992}{2009/08/20}{renewpagestyle里面前导的空格可能导致clearpage生成新的一页，将空格去掉}
% \changes{v0.993}{2009/08/26}{修改标题，博士硕士对应不同的页眉}
%
%    \begin{macrocode}
%<*cls>
\renewpagestyle{plain}{
\sethead{}{\raisebox{.65\baselineskip}{\songti \wuhao \ifisanon{~}\else{国防科学技术大学研究生院\@optionpaperclass{}学位论文}\fi}}{}%
\setfoot{}{{\songti \wuhao 第~\thepage~页}}{}%
\headrule%
\footrule%
}
\setfootrule{1bp}
%</cls>
%    \end{macrocode}
%
%\subsection{编写格式}
%
% 当页面设置好之后，就是在论文的不同部分分别调用，一般来说论文类的书籍
% 分为三个matter，为前言区（前置部分），正文区（主体），后文区（附录），
% 在国防科大论文书写要求中，
% 需要将摘要单独进行页码编号，其编号为小写罗马字母，为此，
% 可以将摘要单独设定为一个matter，
% 名叫就叫做MidMatter，称作摘要区。每个Matter我们都一一介绍。
%
% 首先看前置部分，主要包括封面，目录，摘要等，实现为：
%
%    \begin{macrocode}
%<*cls>
\renewcommand\frontmatter{%
    \if@openright\cleardoublepage\else\clearpage\fi
    \@mainmatterfalse
    \pagenumbering{Roman}
    \pagestyle{plain}}
\newcommand\midmatter{%
    \if@openright\cleardoublepage\else\clearpage\fi
    \@mainmatterfalse
    \pagenumbering{roman}
    \pagestyle{plain}}
%</cls>
%    \end{macrocode}
%
% 之后为文章的正文区，采用阿拉伯数字编页码：
%
%    \begin{macrocode}
%<*cls>
\renewcommand\mainmatter{%
    \if@openright\cleardoublepage\else\clearpage\fi
    \@mainmattertrue
    \pagenumbering{arabic}
    \pagestyle{plain}}
%</cls>
%    \end{macrocode}
%
% 最后是附录部分，由于他的章节标题与正文中不一样(不是第几章，而是附录几)，
% 我们需要单独设定：
%
%    \begin{macrocode}
%<*cls>
\renewcommand\backmatter{%
    \if@openright\cleardoublepage\else\clearpage\fi
    \titleformat{\chapter}{\filcenter \heiti \sanhao}{附录\,\thechapter\,}{1em}{}
    \titlecontents{chapter}[0pt]{\vspace{0.25\baselineskip} \heiti \xiaosi[1.25]}
      {附录\,\thecontentslabel\quad}{}
      {\hspace{.5em}\titlerule*{.}\contentspage}
    \@mainmattertrue
    \pagestyle{plain}}
%</cls>
%    \end{macrocode}
%
% 我们重新定义\verb|cleardoublepage|，使得生成完全的空白页，页面模式为\verb|empty|
%    \begin{macrocode}
%<*cls>
\renewcommand\cleardoublepage{\clearpage\if@openright \ifodd\c@page\else
  \newpage{}
  \thispagestyle{empty}
  \vspace*{\fill}
  \begin{center}
  \end{center}
  \vspace*{\fill}
  \clearpage\fi\fi%
}
%</cls>
%    \end{macrocode}
%
%\subsubsection{前置目录}
% 前置部分的封面在后面详细介绍。首先看目录，要求为：
% 目次页由论文的章、节、条、项、附录等的序号、名称和页码组成，
% 另页排在序之后。目次页标注学位论文的前三级目录。
% 标题统一用“目录”，黑体3字号字居中，段前、段后间距为1行;
% 各章（一级目录）名称用黑体小4号字，段前间距为0.5行，
% 段后间距为0行; 其它（二、三级目录）用宋体小4号字，
% 段前、段后间距为0行。：
%
% 在\LaTeX{}中，Chapter在目录中默认是没有点的,我们加上,另外我们一并将
% 目录中的section和subsection设定好,
% \changes{v0.991}{2009/08/18}{modify TOC baselineskip and font lineskip to 1.25}
%
%    \begin{macrocode}
%<*cls>
\titlecontents{chapter}[0pt]{\vspace{0.25\baselineskip} \heiti \xiaosi[1.25]}
    {第\CJKnumber{\thecontentslabel}章\quad}{}
    {\hspace{.5em}\titlerule*{.}\contentspage}
\titlecontents{section}[2em]{\songti \xiaosi[1.25]}
    {\thecontentslabel\quad}{}
    {\hspace{.5em}\titlerule*{.}\contentspage}
\titlecontents{subsection}[4em]{\songti \xiaosi[1.25]}
    {\thecontentslabel\quad}{}
    {\hspace{.5em}\titlerule*{.}\contentspage}
%</cls>
%    \end{macrocode}
%
% 然后是表目录和图目录，内容用宋体小4号字，在同学使用模板时，需要标题对齐，
% 我们一并在这里实现:
% \changes{v0.993}{2009/08/25}{添加makebox使得图表标题对齐}
%
%    \begin{macrocode}
%<*cls>
\titlecontents{figure}[0pt]{\songti \xiaosi[1.25]}
    {\makebox[3.5em][l]{图~\thecontentslabel\quad}}{}
    {\hspace{.5em}\titlerule*{.}\contentspage}
\titlecontents{table}[0pt]{\songti \xiaosi[1.25]}
    {\makebox[3.5em][l]{表~\thecontentslabel\quad}}{}
    {\hspace{.5em}\titlerule*{.}\contentspage}
%</cls>
%    \end{macrocode}
%
% 书籍模板中，在LOF或者LOT章节之间会默认插入额外的距离，我们通过修改下面这个命令移除，
% 这个方法不是一个完美的办法，\textbf{注意}：下面的代码不要去深究或者理解，
% 这只是把book.cls中的内容复制过来，然后去掉包含addvspace命令的两行。
% 我实在找不出更加好的办法，如果你有，可以联系我。
%
% \changes{v0.993}{2009/08/25}{移除LOF及LOT中章节之间额外的距离}
%
%    \begin{macrocode}
%<*cls>
\renewcommand\chapter{\if@openright\cleardoublepage\else\clearpage\fi
                    \thispagestyle{plain}%
                    \global\@topnum\z@
                    \@afterindentfalse
                    \secdef\nudt@chapter\@schapter}
\def\nudt@chapter[#1]#2{
  \ifnum \c@secnumdepth >\m@ne
    \if@openright\cleardoublepage\else\clearpage\fi
    \phantomsection
    \if@mainmatter
      \refstepcounter{chapter}%
      \addcontentsline{toc}{chapter}%
        {\protect\numberline{\thechapter}#1}%
    \else
      \addcontentsline{toc}{chapter}{#1}%
    \fi
  \else
    \addcontentsline{toc}{chapter}{#1}%
  \fi
  \chaptermark{#1}%
  \if@twocolumn
    \@topnewpage[\@makechapterhead{#2}]%
  \else
    \@makechapterhead{#2}%
    \@afterheading
  \fi
}
%</cls>
%    \end{macrocode}
%
%\subsubsection{前置摘要}
%
% 摘要的要求为题目黑体3字号字居中，段前、段后间距为1行，内容用宋体小4号字，
% 英文摘要内容用Time New Roman小4号字。
% 中文关键字以黑体小4号字另起一行，排在摘要的下方，英文关键字用Arial小4号字。
%
% \changes{v1.8}{2010/10/15}{ABSTRACT和英文关键字需要用Arial字体}
%    \begin{macrocode}
%<*cls>
\newcommand\cabstractname{摘\hspace{1em}要}
\newcommand\eabstractname{ABSTRACT}
\newcommand\ckeywordsname{关键词}
\newcommand\ckeywords[1]{{\hei\xiaosi \ckeywordsname: #1}}
\newcommand\ekeywordsname{Key Words}
\newcommand\ekeywords[1]{\textsf{\xiaosi \ekeywordsname: #1}}
\newenvironment{cabstract}{%
    \chapter{\cabstractname}
    \xiaosi
    \@afterheading}
    {\par\vspace{2em}\par}
\newenvironment{eabstract}{%
    \chapter{\textsf{\eabstractname}}
    \xiaosi
    \@afterheading}
    {\par\vspace{2em}\par}
%</cls>
%    \end{macrocode}
%
%\subsection{主体部分}
%
% \subsubsection{标题格式}
% 要求为：
% \begin{compactenum}
% \item	一级标题(章)用黑体3号字居中，1.25倍行距，段前、段后间距为1行，每一章从新的一页开始;
% \item	二级标题(节)用宋体4号粗体字居中，1.25倍行距，段前、段后间距为1行;
% \item	三级标题用黑体小4号字两端对齐，1.25倍行距，段前、段后间距为1行;
% \item	四级标题用宋体小4号粗体字两端对齐，1.25倍行距，段前间距为0.5行，段后间距为0行;
% \end{compactenum}
%
% \changes{v0.991}{2009/08/18}{按照要求设定标题}
% \changes{v0.992}{2009/08/19}{修改secnumdepth使得subsubsection可用}
% \changes{v1.1}{2009/09/26}{修改Title的spacing为弹性值}
% \changes{v1.2}{2009/10/06}{去掉弹性值，不去生成大量的空白}
% \changes{v1.4}{2009/10/28}{修改chapter段后行距为2ex，段前-1ex，保证上下对称}
% \changes{v1.4}{2009/10/29}{修改chapter段后行距为2.4ex，段前-1.2ex，保证上下对称}
%
% 当章节标题出现的新的一页时，会出现段前距过小的情况，按照milksea的说法是：
% 一般而言，当一个内容在一页开头时，前面的\verb|\vskip|不起作用；
% 类似地，一行开头\verb|\hskip|不起作用。这不是 BUG，如果需要总起效果的间距，
% 用\verb|\vspace*|，文档里面有这样的例子。参照titlesec的文档，需加上:
% \changes{v1.9}{2010/10/28}{增加sectionbreak，设定topskip为0pt}
%
%    \begin{macrocode}
%<*cls>
\newcommand{\sectionbreak}{%
\addpenalty{-300}%
\vspace*{0pt}%
}
\setlength{\topskip}{0pt}
%</cls>
%    \end{macrocode}
% \changes{v1.9}{2010/10/28}{在定义了粗宋字体之后，按照学位论文要求设定标题字体}
% \changes{v1.9}{2010/10/28}{使用了ttltips.pdf的设置chapter距顶端距离的办法}
% \changes{v2.3}{2013/12/27}{修改了footnotesize，当一页中有两个footnote时，间距调整正常}
% \changes{v2.3}{2013/12/27}{修改了titleformat和titlespacing的定制，更改英文字体，修正间距}
%
%    \begin{macrocode}
%<*cls>
\setcounter{secnumdepth}{3}
\setlength{\footnotesep}{2.2ex \@minus 2bp}
\titleformat{\chapter}{\filcenter\sf \heiti\sanhao[1.25]}{第\CJKnumber{\thechapter}章\,}{1em}{}
\titleformat{\section}{\filcenter\bfseries \cusong\sihao[1.25]}{\thesection}{1em}{}
\titleformat{\subsection}{\sf \heiti\xiaosi[1.25]}{\thesubsection}{1em}{}
\titleformat{\subsubsection}{\bfseries \cusong\xiaosi[1.25]}{\thesubsubsection}{1em}{}
\titlespacing{\chapter}{0pt}{2.4ex-\topskip-\heightof{A}}{2.4ex \@minus 2bp}
\titlespacing{\section}{0pt}{2ex-\heightof{a}}{2ex \@minus 2bp}
\titlespacing{\subsection}{2em}{2ex \@minus 2bp}{2ex \@minus 2bp}
\titlespacing{\subsubsection}{2em}{1ex \@minus 2bp}{0ex \@minus 2bp}
%</cls>
%    \end{macrocode}
%
%\subsubsection{正文字体}
% 首先确定正文中使用的字体，文档要求正文字体为小四，行距为固定值1.25倍，
% 中文字体为宋体，英文为{Times New Roman}
%
%\begin{macro}{\normalsize}
% 我们重新定义 |\normalsize| 来确定文档的正文字体，
% 同时修改正文中公式与文字间的距离：
% \changes{v1.9}{2010/10/28}{在normalsize后面每一行加上\%号来吃掉多余的空格}
% \changes{v2.2}{2011/07/16}{减小公式之间距离rubber space的上界}
% \changes{v2.2}{2011/09/25}{减小公式之间距离rubber space的下界}
% \changes{v2.3}{2013/12/27}{移除了abovedisplay的正向rubberspace}
%
%    \begin{macrocode}
%<*cls>
\renewcommand\normalsize{%
\@setfontsize\normalsize{12bp}{12.87bp}%
\renewcommand{\baselinestretch}{1.3}%
\setlength\abovedisplayskip{10bp \@minus 1bp}%
\setlength\abovedisplayshortskip{10bp \@minus 1bp}%
\setlength\belowdisplayskip{\abovedisplayskip}%
\setlength\belowdisplayshortskip{\abovedisplayshortskip}%
}
%</cls>
%    \end{macrocode}
%\end{macro}
%
% \changes{v0.991}{2009/08/18}{modify normalsize, which will cause headrule shift}
% \changes{v0.991}{2009/08/18}{add comment on displayskip}
% \changes{v1.0}{2009/09/22}{modify display skip}
%
%\subsubsection{正文段落}
% 接下来还有一个细节就是处理段落缩进，文档设定为首行缩进2个字符，
% 这一个命令需要在文档开始时自动执行：
%
% \changes{v1.3}{2009/10/03}{添加checkparameter这一选项，避免由于更新模板导致未定义的情况出现}
% \changes{v1.7}{2010/04/30}{应当在封面制作完后替换tabular}
%
%    \begin{macrocode}
%<*cls>
\newlength\CJK@twochars
\def\CJK@spaceChar{\Unicode{48}{7}}
\def\CJKindent{%
  \settowidth\CJK@twochars{中国}%
  \parindent\CJK@twochars}
\AtBeginDocument{%
  \CJKindent\relax
  \checkparameter\relax
}
%</cls>
%    \end{macrocode}
%
% 之后定义段落间距，段前间距以及段后间距都为0
% \changes{v0.993}{2009/08/27}{修改parskip}
% \changes{v2.2}{2011/09/25}{修改parskip，允许少量的调整，1bp}
% \changes{v2.2}{2011/10/14}{修改parskip，仅允许负的少量调整，2bp}
% \changes{v2.3}{2013/12/27}{移除parskip的正向rubberspace}
%
%    \begin{macrocode}
%<*cls>
\setlength{\parskip}{0bp \@minus 2bp}
%</cls>
%    \end{macrocode}
%
% 有时候我们需要手动设定字体间距，该命令在声明页使用过：
%\begin{macro}{\ziju}
%    \begin{macrocode}
%<*cls>
\newcommand*{\ziju}[1]{\renewcommand{\CJKglue}{\hskip #1}}
%</cls>
%    \end{macrocode}
%\end{macro}
%
% \changes{v1.4}{2009/10/26}{推荐用户使用紧凑的列表环境}
%
% 这一部分来自Thuthesis的代码，其出发点是不满意\LaTeX{}默认列表环境间距过大，用
% paralist包中的相关环境进行替代。请参考paralist宏包。
%
% \changes{v1.4}{2009/10/26}{修改参考文献的行距设定}
%
% 而同样有间距问题的是参考文献，两个条目之间过大的距离不是很美观，
% 最简单的办法是修改bibsep变量，如果还是不行，我们直接从thuthesis中拿来代码:
%
% \changes{v1.4}{2009/10/26}{修改参考文献的行距}
% \changes{v1.4}{2009/10/29}{修改参考文献左对齐}
% \changes{v2.2}{2011/10/14}{减小文献列表间距，将penalty改为4000}
%
%    \begin{macrocode}
%<*cls>
\renewenvironment{thebibliography}[1]{%
   \chapter*{\bibname}%
   \list{\@biblabel{\@arabic\c@enumiv}}%
        {\renewcommand{\makelabel}[1]{##1\hfill}
         \settowidth\labelwidth{1.1cm}
         \setlength{\labelsep}{0.4em}
         \setlength{\itemindent}{0pt}
         \setlength{\leftmargin}{\labelwidth+\labelsep}
         \addtolength{\itemsep}{-0.7em}
         \usecounter{enumiv}%
         \let\p@enumiv\@empty
         \renewcommand\theenumiv{\@arabic\c@enumiv}}%
    \sloppy\frenchspacing
    \clubpenalty4000%
    \widowpenalty4000%
    \interlinepenalty4000%
    \sfcode`\.\@m}
   {\def\@noitemerr
     {\@latex@warning{Empty `thebibliography' environment}}%
    \endlist\frenchspacing}
%</cls>
%    \end{macrocode}
%
%\subsection{浮动对象}
%
% 浮动对象针对的目标是图片表格，标题为五号字体，
% 图片标题在下，表格标题在上，具体实现为:
% \changes{v1.1}{2009/09/26}{修改float浮动弹性}
% \changes{v2.2}{2011/09/25}{去掉1.0fil改为4bp, 这样不至于生成过大的空白.}
% \changes{v2.3}{2013/12/27}{移除了floatsep所有rubberspace. 原先为正2负2.}
%
%    \begin{macrocode}
%<*cls>
\setlength{\floatsep}{12bp}
\setlength{\intextsep}{12bp}
\setlength{\textfloatsep}{12bp}
\setlength{\@fptop}{0bp}
\setlength{\@fpsep}{12bp}
\setlength{\@fpbot}{0bp}
%</cls>
%    \end{macrocode}
%
% 接下来设置每一页图形占据的比例，这个直接从\thuthesis{}中拿出，
% 具体含义可以参考下面这个网页：
% \url{http://www.ctex.org/documents/latex/graphics/node69.html}，
% 里面解释的很清楚，这个布置方法也是网站的推荐：
% \changes{v1.3}{2009/09/29}{调整floatpagefraction的大小}
% \changes{v2.2}{2011/09/10}{重新调整floatpagefraction，使得更为宽松}
% \changes{v2.2}{2011/09/25}{更为宽松的布置条件}
% \changes{v2.2}{2011/09/25}{更为宽松的布置条件，仿照AMSMATH}
%
%    \begin{macrocode}
%<*cls>
\renewcommand{\textfraction}{0.01}
\renewcommand{\topfraction}{0.99}
\renewcommand{\bottomfraction}{0.99}
\renewcommand{\floatpagefraction}{0.90}
\clubpenalty            =   10000
\widowpenalty           =   10000
\displaywidowpenalty    =   10000
%</cls>
%    \end{macrocode}
%
% 在修改图片标题距离时，要注意，aboveskip为内距离，也就是标题与浮动体之间的距离，
% belowskip为外距离，也就是标题与正文之间的距离。
% \changes{v1.3}{2009/09/29}{缩小图片标题与下文的距离}
% \changes{v1.7}{2010/04/30}{增添LT array命令，可修改Longtable字体大小}
%
%    \begin{macrocode}
%<*cls>
\let\old@tabular\@tabular
\def\thu@tabular{\wuhao[1.25]\old@tabular}
\DeclareCaptionLabelFormat{thu}{{\wuhao[1.25]\song #1~\rmfamily #2}}
\DeclareCaptionLabelSeparator{thu}{\hspace{1em}}
\DeclareCaptionFont{thu}{\wuhao[1.25]}
\captionsetup{labelformat=thu,labelsep=thu,font=thu}
\captionsetup[table]{position=top,belowskip=0bp \@plus 2bp \@minus 2bp,aboveskip=6bp \@plus 2bp \@minus 2bp}%
\captionsetup[figure]{position=bottom,belowskip=-3bp \@plus 2bp \@minus 2bp,aboveskip=6bp \@plus 2bp \@minus 2bp}%
\captionsetup[subfloat]
{labelformat=simple,font=thu,captionskip=6bp,nearskip=6bp,farskip=0bp,topadjust=0bp}
\renewcommand{\thesubfigure}{(\alph{subfigure})}
\renewcommand{\thesubtable}{(\alph{subtable})}
\let\thu@LT@array\LT@array
\def\LT@array{\thu@LT@array}
%</cls>
%    \end{macrocode}
%
%\subsection{自定环境}
%
% 在这里我们自定义一些论文种会使用到的环境，主要有摘要，符号表，致谢，个人介绍等：
% 这些单独定义的环境可以分别配置以满足要求。
%
% 有些论文需要在正文前面加入符号列表, 其内容格式是简单的列表环境:
% \changes{v2.2}{2011/09/10}{略微减小列表间距，与行距相等}
% \changes{v2.2}{2011/09/13}{略微减小标签和说明间距}
% \changes{v2.3}{2013/12/27}{重新制作了denotation, 参见样例文件}
%
%    \begin{macrocode}
%<*cls>
\newenvironment{denotation}[1][2.5cm]{
    \noindent\begin{list}{}%
    {\vskip-30bp\xiaosi[1.5]
    \renewcommand\makelabel[1]{##1\hfil}
    \setlength{\labelwidth}{#1} % 标签盒子宽度
    \setlength{\labelsep}{0cm} % 标签与列表文本距离
    \setlength{\itemindent}{2em} % 标签缩进量
    \setlength{\leftmargin}{\labelwidth+\labelsep} % 左边界
    \setlength{\rightmargin}{0cm}
    \setlength{\parsep}{0cm} % 段落间距
    \setlength{\itemsep}{0cm} % 标签间距
    \setlength{\listparindent}{0cm} % 段落缩进量
    \setlength{\topsep}{0pt} % 标签与上文的间距
}}{\end{list}}
%</cls>
%    \end{macrocode}
%
% 致谢往往在正文的最后：
%
% \changes{v1.8}{2010/10/15}{致谢之间要加一个空格}
%    \begin{macrocode}
%<*cls>
\newenvironment{ack}{%
    \chapter*{致\hspace{1em}谢}%
    \addcontentsline{toc}{chapter}{致谢}%
    \ifisanon\color{white}\else\relax\fi%
    \xiaosi%
    \@afterheading}
    {\par\vspace{2em}\par}
%</cls>
%    \end{macrocode}
%
% 个人简历这一部分用来放置作者在研究生期间取得的成果，发表的论文等。可以
% 详细的参考\verb|data/|中的文件自己书写。
% \changes{v1.4}{2009/10/26}{修改标题}
%
%    \begin{macrocode}
%<*cls>
\newenvironment{resume}{%
    \chapter*{作者在学期间取得的学术成果}
    \addcontentsline{toc}{chapter}{作者在学期间取得的学术成果}
    \xiaosi
    \@afterheading}
    {\par\vspace{2em}\par}
%</cls>
%    \end{macrocode}
%
%\subsubsection{定理环境}
% 定理环境可能数学论文中应用较多:
% \changes{v2.0}{2010/11/09}{修改定理的分隔符和QED符号，修改字体;缩进为段落缩进;修改编号}
% \changes{v2.2}{2011/10/14}{设定定理、定义环境的合理间隔}
%
%    \begin{macrocode}
%<*cls>
\renewtheoremstyle{nonumberplain}%
{\item[\hspace*{2em} \theorem@headerfont ##1\ \theorem@separator]}%
{\item[\hspace*{2em} \theorem@headerfont ##1\ (##3)\theorem@separator]}
\theoremstyle{nonumberplain}
\theorembodyfont{\kai\xiaosi[1.3]}
\theoremheaderfont{\hei\xiaosi[1.3]}
\theoremsymbol{\ensuremath{\blacksquare}}
\theoremseparator{:\,}
\newtheorem{proof}{证明}[chapter]
\newtheorem{assumption}{假设}[chapter]
\newtheorem{definition}{定义}[chapter]

\renewtheoremstyle{plain}%
{\item[\hspace*{2em} \theorem@headerfont ##1\ ##2\theorem@separator]}%
{\item[\hspace*{2em} \theorem@headerfont ##1\ ##2\ (##3)\theorem@separator]}
\theoremstyle{plain}
\theorembodyfont{\kai\xiaosi[1.3]}
\theoremheaderfont{\hei\xiaosi[1.3]}
\theoremsymbol{}
\newtheorem{lemma}{引理}[chapter]
\newtheorem{theorem}{定理}[chapter]
\newtheorem{axiom}{公理}[chapter]
\newtheorem{corollary}{推论}[chapter]
\newtheorem{conjecture}{猜想}[chapter]
\newtheorem{proposition}{命题}[chapter]
\newtheorem{exercise}{练习}[section]
\newtheorem{example}{例}[section]
\newtheorem{problem}{问题}[section]
\newtheorem{remark}{注释}[section]
%</cls>
%    \end{macrocode}
%
%\subsection{论文属性}
% 这里的内容主要用来定义封面中的一些元素，你可以像填空一样完成封面的制作：
% \changes{v0.993}{2009/08/26}{添加cosupervisor，协助指导导师}
% \changes{v1.3}{2009/10/03}{添加英文第二导师，增加判断语句，在文档开始时执行}
%
%    \begin{macrocode}
%<*cls>
\def\classification#1{\def\@classification{#1}} % 中图分类号
\def\serialno#1{\def\@serialno{#1}} % 学号
\def\UDC#1{\def\@UDC{#1}} % UDC号
\def\confidentiality#1{\def\@confidentiality{#1}} % 密级
\def\title#1{\def\@title{#1}} % 中文题目
\newtoks\displaytitle
\def\author#1{\def\@author{#1}}
\def\zhdate#1{\def\@zhdate{#1}}	% 中文日期
\def\subject#1{\def\@subject{#1}} % 中文学科
\def\researchfield#1{\def\@researchfield{#1}} % 中文研究方向
\def\supervisor#1{\def\@supervisor{#1}} % 导师
\def\cosupervisor#1{\def\@cosupervisor{#1}} % 协助指导教师
\def\papertype#1{\def\@papertype{#1}} % 工学，理学，同等学历申请工(理)学
\def\entitle#1{\def\@entitle{#1}}
\def\enauthor#1{\def\@enauthor{#1}}
\def\ensupervisor#1{\def\@ensupervisor{#1}}
\def\encosupervisor#1{\def\@encosupervisor{#1}}
\def\endate#1{\def\@endate{#1}}
\def\ensubject#1{\def\@ensubject{#1}}
\def\enpapertype#1{\def\@enpapertype{#1}} % Engineering, Science
\def\optionpaperclass#1{\def\@optionpaperclass{#1}} % paperclass
\def\optionpaperclassen#1{\def\@optionpaperclassen{#1}}
\def\optionas#1{\def\@optionas{#1}} % Advisor OR Supervisor
%</cls>
%    \end{macrocode}
%
% 我们看用户是想用博士封面还是硕士封面:
%
%    \begin{macrocode}
%<*cls>
\ifismaster
  \optionpaperclass{硕士}
  \optionpaperclassen{Master}
  \optionas{Advisor}
\else
  \optionpaperclass{博士}
  \optionpaperclassen{Doctor}
  \optionas{Supervisor}
\fi
%</cls>
%    \end{macrocode}
%
%\subsection{制作封面}
%
% 由于封面中一些元素是可选的，如果在正文中没有定义，那么判断ifx的时候就会出错，
% 我们加入下面的命令进行判断，如果没定义，我们就令他为空。
% 这个命令将在文档开始时自动执行。
%
%    \begin{macrocode}
%<*cls>
\newcommand{\checkparameter}
{
  \ifthenelse{\isundefined{\@cosupervisor}}{\cosupervisor{}}{}
  \ifthenelse{\isundefined{\@encosupervisor}}{\encosupervisor{}}{}
}
%</cls>
%    \end{macrocode}
%
% 制作封面比较复杂，需要一些手动调整的东西，首先来看第一页，
% 重新定义了\verb|maketitle|，
% 用表格来安排页面元素，页头采用仿宋五号字体，段前段后间距一行，这个空一行就用3ex实现，
% \changes{v1.5}{2009/11/18}{微调封面布局}
% \changes{v2.0}{2010/11/09}{标题用粗宋}
% \changes{v2.0}{2010/11/09}{增加盲评控制}
% \changes{v2.1}{2010/11/23}{设定页码为alph，增加cleardoublepage，这样在博士双面时方便打印}
%    \begin{macrocode}
%<*cls>
\def\maketitle{%
  \def\entry##1##2##3{%
    \multicolumn{##1}{l}{\underline{\hbox to ##2{\hfil##3\hfil}}}
    }
  \null
  \ifisanon%
  \author{}%
  \enauthor{}%
  \supervisor{}%
  \cosupervisor{}%
  \ensupervisor{}%
  \encosupervisor{}%
  \else\relax\fi%
  \pagenumbering{alph}% not display, for print only
  \thispagestyle{empty}%
  \begin{center}\leavevmode	% 表格环境
  {\fangsong \wuhao[1.25]%
    \begin{tabular}{llcll}
    分类号 	& \entry{1}{3.2cm}{\@classification} & \hspace*{4.8cm}%
    学号   	& \entry{1}{3.2cm}{\@serialno}         \\[5mm]    %
    U\ D\ C	& \entry{1}{3.2cm}{\@UDC} &            \hspace*{4.8cm}
    密级	& \entry{1}{3.2cm}{\@confidentiality}
    \end{tabular}
  }
  \par
  \vspace*{2.5cm} %插入空白
  {\heiti\sanhao \@papertype{}\@optionpaperclass{}学位论文}\\
  \vspace{12bp}
  {\cusong\erhao[1.25] \@title \par}%
  \vspace{45bp} %从WORD中得来
  {\heiti \sihao
    \begin{tabular}{cp{8cm}c}
      \raisebox{-3.7ex}[0pt]{\@optionpaperclass{}生姓名} &
        {\fs \hfil\raisebox{-3.7ex}[0pt]{\@author}\hfil{}} & \\[3.2ex]
        \cline{2-2}
      \raisebox{-3.7ex}[0pt]{学\ 科\ 专\ 业} &
        {\fs \hfil\raisebox{-3.7ex}[0pt]{\@subject}\hfil{}} & \\[3.2ex]
        \cline{2-2}
      \raisebox{-3.7ex}[0pt]{研\ 究\ 方\ 向} &
        {\fs \hfil\raisebox{-3.7ex}[0pt]{\@researchfield}\hfil{}} & \\[3.2ex]
        \cline{2-2}
      \raisebox{-3.7ex}[0pt]{指\ 导\ 教\ 师} &
        {\fs \hfil\raisebox{-3.7ex}[0pt]{\@supervisor}\hfil{}} & \\[3.2ex]
        \cline{2-2}
      \ifx\@cosupervisor\@empty\else
        & {\fs \hfil\raisebox{-3.7ex}[0pt]{\@cosupervisor}\hfil{}} & \\[3.2ex]
        \cline{2-2}
      \fi
    \end{tabular}
  }
  \end{center}%

  \par
  \vfill
  {\centering \cusong \sanhao \ifisanon{~}\else{国防科学技术大学研究生院}\fi\\[0.8em]
    {\@zhdate \par}%
  }
  \vspace{1mm}
%</cls>
%    \end{macrocode}
%
%第二页主要是论文的英文信息，简称英文封面
%
%    \begin{macrocode}
%<*cls>
  \cleardoublepage%
  \newpage
  \thispagestyle{empty}%

  \begin{center}\leavevmode
  \vfill\bfseries
  {\erhao[1.25] \@entitle \par}
  {\sanhao[1.25]
  \vfill\vfill\vfill\vfill\vfill\vfill
  \begin{tabular}{rl}
    Candidate:\ & {\textsf{\@enauthor}}\\
    \@optionas{}:\ & {\textsf{\@ensupervisor}}\\
    \ifx\@encosupervisor\@empty\else
      & {\textsf{\@encosupervisor}} \\
    \fi
  \end{tabular}}
  \vfill\vfill\vfill\vfill
  {\sanhao[1.5]
A dissertation\\
Submitted in partial fulfillment of the requirements\\
for the degree of \textsf{\@optionpaperclassen{} of \@enpapertype}\\
in \textsf{\@ensubject}\\
\makebox[\textwidth]{\ifisanon{~}\else{Graduate School of National University of %
Defense Technology}\fi}\\
\ifisanon{~}\else{Changsha, Hunan, P.\ R.\ China}\fi\\[5mm]
~\@endate~
  }
  \end{center}\vfill
  \cleardoublepage%
%</cls>
%    \end{macrocode}
%
% 第三页放置独创性声明，这里要使用\verb|displaytitle|这个论文元素：
% \changes{v2.1}{2010/12/29}{表格字体定为五号}
%    \begin{macrocode}
%<*cls>
  \newpage
  \thispagestyle{empty}

  {\cusong \erhao \centering \ziju{12pt} 独创性声明 \par\vspace{2cm}}
    \renewcommand{\baselinestretch}{1.5}%
  {\fangsong\xiaosi %
本人声明所呈交的学位论文是我本人在导师指导下进行的研究工
作及取得的研究成果。尽我所知，除文中特别加以标注和致谢的地方外，论文中
不包含其他人已经发表和撰写过的研究成果，也不包含为获得国防科学技术大学或
其他教育机构的学位或证书而使用过的材料。与我一同工作的同志对本研究所做的
任何贡献均已在论文中作了明确的说明并表示谢意。\par
学位论文题目：\vbox{\hbox to11cm{\hfil \the\displaytitle \hfil}
  \protect\vspace{0.6truemm}\relax
  \hrule depth0pt height0.15truemm width11cm}\par
学位论文作者签名：\hrulefill\hrulefill\hrulefill\hrulefill\hrulefill
  \hfill 日期：\hfill\hfill 年\hfill 月 \hfill 日\hspace{1cm}\par}

  \vspace*{2cm}
  {\cusong \erhao \centering 学位论文版权使用授权书\par\vspace{2cm}}
  {\fangsong\xiaosi %
本人完全了解国防科学技术大学有关保留、使用学位论文的规定。
本人授权国防科学技术大学可以保留并向国家有关部门或机构送交论文的复印件和电子
文档，允许论文被查阅和借阅; 可以将学位论文的全部或部分内容编入有关数据库
进行检索，可以采用影印、缩印或扫描等复制手段保存、汇编学位论文。\par
（保密学位论文在解密后适用本授权书。）\par
学位论文题目：\vbox{\hbox to11cm{\hfil \the\displaytitle \hfil}
  \protect\vspace{0.6truemm}\relax
  \hrule depth0pt height0.15truemm width11cm}\par
学位论文作者签名：\hrulefill\hrulefill\hrulefill\hrulefill\hrulefill
  \hfill 日期：\hfill\hfill 年\hfill 月 \hfill 日\par
作者指导教师签名：\hrulefill\hrulefill\hrulefill\hrulefill\hrulefill
  \hfill 日期：\hfill\hfill 年\hfill 月 \hfill 日\par}

  \normalsize % normal, 正文开始
  \def\@tabular{\wuhao[1.25]\old@tabular} % 之后表格字体使用5号
  \cleardoublepage%
  \newpage
  \thispagestyle{empty}

} % QED
%</cls>
%    \end{macrocode}
%
% \Finale
%
% \iffalse meta-comment
% 下面的配置为本文档使用的模板
% \fi
%
% \iffalse
% \changes{v1.8}{2010/10/15}{按照milksea的说法，修改xltxtra和xunicode相对于xeCJK的顺序}
% \changes{v1.8}{2010/10/15}{添加了xeCJKsetcharclass，使itemize环境正确}
%<*nudtx>
%    \begin{macrocode}
%    \end{macrocode}
%</nudtx>
% \fi
%
\endinput
